% =============================================================
% Terimler Sözlüğü -- Türkçe-İngilizce / İngilizce-Türkçe
% Bilgisayar Mimarisi Kitabı
% Bu dosya % =============================================================
% Terimler Sözlüğü -- Türkçe-İngilizce / İngilizce-Türkçe
% Bilgisayar Mimarisi Kitabı
% Bu dosya % =============================================================
% Terimler Sözlüğü -- Türkçe-İngilizce / İngilizce-Türkçe
% Bilgisayar Mimarisi Kitabı
% Bu dosya % =============================================================
% Terimler Sözlüğü -- Türkçe-İngilizce / İngilizce-Türkçe
% Bilgisayar Mimarisi Kitabı
% Bu dosya \input{terimler-sozlugu} ile dahil edilir.
% \chapter komutu ana-kitap.tex dosyasında tanımlanmalıdır.
% =============================================================

%% ============================================================
%%  BÖLÜM 1: Türkçe → İngilizce
%%  Türkçe alfabetik sıra: a,b,c,ç,d,e,f,g,ğ,h,ı,i,j,k,l,m,
%%                          n,o,ö,p,r,s,ş,t,u,ü,v,y,z
%% ============================================================
\section*{Türkçe -- İngilizce}
\addcontentsline{toc}{section}{Türkçe -- İngilizce}

{\small
\begin{longtable}{@{}p{7.6cm}p{7.6cm}@{}}
\toprule
\textbf{Türkçe Terim} & \textbf{İngilizce Karşılık} \\
\midrule
\endfirsthead

\toprule
\textbf{Türkçe Terim (devam)} & \textbf{İngilizce Karşılık} \\
\midrule
\endhead

\midrule
\multicolumn{2}{r}{\small\textit{(devamı sonraki sayfada)}} \\
\endfoot

\bottomrule
\endlastfoot

%% --- A ---
adresleme kipleri & addressing modes \\
ağırlıklı ortalama & weighted average \\
Akış Çoklu İşlemci (\texttt{SM}) & Streaming Multiprocessor (\texttt{SM}) \\
altyordam & subroutine \\
Amdahl Yasası & Amdahl's Law \\
\texttt{AMB} & \texttt{ALU} (Arithmetic Logic Unit) \\
anlamlı kısım & significand / mantissa \\
anlık değer & immediate value \\
anlık değer üreteci & immediate generator \\
ara bellek & buffer \\
ardışık tutarlılık & sequential consistency \\
arama zamanı & seek time \\
aritmetik buyruklar & arithmetic instructions \\
aşınma dengeleme & wear leveling \\
aritmetik mantık birimi (\texttt{AMB}) & Arithmetic Logic Unit (\texttt{ALU}) \\
aritmetik yoğunluk & arithmetic intensity \\
artık-N gösterimi & excess-N notation \\
\texttt{ASCII} & \texttt{ASCII} \\
\texttt{ASIC} & Application-Specific Integrated Circuit \\
çevirici dil (Assembly) & assembly language \\
atlama buyrukları & jump instructions \\
atomik işlem & atomic operation \\
aygıt & device \\
ayrıcalıklı kip & privileged mode \\

%% --- B ---
bilimsel gösterim & scientific notation \\
Booth algoritması & Booth's algorithm \\
B türü & B-type \\
bağlam değiştirme & context switching \\
bağlantı noktası & port \\
bağlayıcı & linker \\
bakış tablosu (\texttt{LUT}) & Look-Up Table (\texttt{LUT}) \\
bant genişliği & bandwidth \\
başarım & performance \\
bayat veri & stale data \\
bayt & byte \\
bayt sıralaması & byte order \\
büyük sayfa & huge page \\
\texttt{BBÇ} & \texttt{CPI} (cycles per instruction) \\
\texttt{BCD} & Binary-Coded Decimal \\
bellek & memory \\
bellek birleştirme & memory coalescing \\
bellek denetleyicisi & memory controller \\
bellek duvarı & memory wall \\
bellek koruması & memory protection \\
bellek hizalama & memory alignment \\
bellek hiyerarşisi & memory hierarchy \\
bellek tutarlığı & cache coherence \\
bellek tutarlık modeli & memory consistency model \\
bellekle eşlemeli G/Ç & memory-mapped I/O \\
benchmark & sınama programı (denektaşı) \\
big.LITTLE & big.LITTLE \\
bire tümleyen & one's complement \\
birleşik bellek mimarisi & unified memory architecture \\
birleşik devre & combinational circuit \\
bit & bit \\
blok & block \\
boru hattı & pipeline \\
boru hattı aşaması & pipeline stage \\
boru hattı kaydı & pipeline register \\
boru hattı tehlikesi & pipeline hazard \\
boşluk (boru hattı) & bubble / NOP \\
bölme & division \\
bulamama (önbellekte) & miss (cache miss) \\
bulamama cezası & miss penalty \\
bulamama oranı & miss rate \\
bulma (önbellekte) & hit (cache hit) \\
bulma oranı & hit rate \\
buyruk & instruction \\
buyruk başına çevrim (\texttt{BBÇ}) & cycles per instruction (\texttt{CPI}) \\
buyruk belleği & instruction memory \\
buyruk çözme & instruction decode \\
buyruk düzeyi paralelliği & instruction-level parallelism (\texttt{ILP}) \\
buyruk biçimi & instruction format \\
buyruk getirme birimi & instruction fetch unit \\
buyruk kümesi mimarisi & Instruction Set Architecture (\texttt{ISA}) \\
buyruk önbelleği & instruction cache \\
buyruk sayacı & instruction counter \\
büyüğü başta & big-endian \\

%% --- C ---
chiplet & chiplet \\
\texttt{CISC} & Complex Instruction Set Computer \\
\texttt{CMOS} & Complementary \texttt{MOS} \\
\texttt{CUDA} & Compute Unified Device Architecture \\

%% --- Ç ---
\texttt{ÇBB} & \texttt{IPC} (instructions per cycle) \\
çağrı kuralları & calling conventions \\
çarpma & multiplication \\
çarpma-toplama (\texttt{MAC}) & multiply-accumulate (\texttt{MAC}) \\
çatı çizgisi modeli & roofline model \\
çekirdek & core \\
çekirdek fonksiyonu (GPU) & kernel function \\
çerçeve göstergesi & frame pointer \\
çevirici & assembler \\
çevrim & cycle \\
çevrim başına buyruk (\texttt{ÇBB}) & instructions per cycle (\texttt{IPC}) \\
çevrim zamanı & cycle time \\
çıkarım (yapay zeka) & inference \\
çıkarma & subtraction \\
yonga içi bağlantı ağı (\texttt{NoC}) & Network on Chip (\texttt{NoC}) \\
çok çekirdekli işlemci & multi-core processor \\
çok seviyeli sayfa tablosu & multi-level page table \\
çok çevrimli & multi-cycle \\
çok vuruşluk işlemci & multi-cycle processor \\
çoklayıcı & multiplexer \\
çözgü & warp \\
çözgü sapması & warp divergence \\

%% --- D ---
dağıtık bellek & distributed memory \\
durum yazmacı & state register \\
dalga eldeli toplayıcı & ripple carry adder \\
dallanma buyrukları & branch instructions \\
dallanma hedef ara belleği & branch target buffer (BTB) \\
dallanma öngörüsü & branch prediction \\
darboğaz & bottleneck \\
\texttt{DDR} & Double Data Rate \\
DEĞİL kapısı & \texttt{NOT} gate \\
dekoder & decoder \\
denetim birimi & control unit \\
denetim sinyali & control signal \\
denetim tablosu & control table \\
denetim tehlikesi & control hazard \\
devingen güç tüketimi & dynamic power consumption \\
devingen öngörü & dynamic prediction \\
Dennard ölçeklemesi & Dennard scaling \\
derin öğrenme & deep learning \\
derleyici & compiler \\
dikey mikroprogramlama & vertical microprogramming \\
dizin tabanlı protokol & directory-based protocol \\
\texttt{DMA} & Direct Memory Access \\
\texttt{DRAM} & Dynamic Random Access Memory \\
doğrudan eşlemeli önbellek & direct-mapped cache \\
doğruluk tablosu & truth table \\
dönme gecikmesi & rotational latency \\
doğrudan yazma & write-through \\
doku belleği & texture memory \\
dolanıklık & entanglement \\
doluluk & occupancy \\
donanım & hardware \\
donanımsal çok kanallılık & hardware multithreading \\
dönüş adresi & return address \\
döngüsel kilit & spin-lock \\
durağan öngörü & static prediction \\
durağan güç tüketimi & static power consumption (leakage) \\
durma & stall \\
düzenleme birimi & hazard detection unit \\

%% --- E ---
\texttt{EDVAC} & \texttt{EDVAC} \\
eğitim (yapay zeka) & training \\
elde & carry \\
elden önceden üretme & carry lookahead \\
\texttt{ENIAC} & \texttt{ENIAC} \\
engel & barrier \\
eşlik (RAID) & parity (RAID) \\
erişim süresi & access time \\
eşzamanlama & synchronization \\
eşzamanlı çok kanallılık (\texttt{SMT}) & Simultaneous Multithreading (\texttt{SMT}) \\
eşzamanlılık & concurrency \\
etiket & tag \\
etkin adres & effective address \\

%% --- F ---
Flaş çeviri katmanı (\texttt{FTL}) & Flash Translation Layer (\texttt{FTL}) \\
flaş bellek & flash memory \\
flip-flop & flip-flop \\
\texttt{FPGA} & Field-Programmable Gate Array \\

%% --- G ---
gecikme & latency \\
geometrik ortalama & geometric mean \\
gecikme odaklı tasarım & latency-oriented design \\
geciktirilmiş dallanma & delayed branching \\
gerçek veri bağımlılığı & true data dependence \\
genel amaçlı yazmaç & general-purpose register \\
genel bellek & global memory \\
geri yazma & write-back \\
getir-çöz-yürüt döngüsü & fetch-decode-execute cycle \\
gevşek tutarlılık & relaxed consistency \\
giriş/çıkış (\texttt{G/Ç}) & input/output (\texttt{I/O}) \\
gölgelendirici & shader \\
gözetleme protokolü & snooping protocol \\
\texttt{GPGPU} & General-Purpose computing on GPU \\
\texttt{GPU} & Graphics Processing Unit \\
gshare öngörücü & gshare predictor \\
Gustafson Yasası & Gustafson's Law \\
güç duvarı & power wall \\
güç tüketimi & power consumption \\

%% --- Ğ ---

%% --- H ---
Harvard mimarisi & Harvard architecture \\
\texttt{HBM} & High Bandwidth Memory \\
heterojen hesaplama & heterogeneous computing \\

%% --- Önbellekte bulamama terimleri (B harfi altında da var) ---
%% --- I ---
ızgara & grid \\

%% --- İ ---
I türü & I-type \\
\texttt{IEEE 754} & \texttt{IEEE 754} \\
yol & bus \\
ikincil bellek & secondary storage \\
\texttt{IOPS} & Input/Output Operations Per Second \\
ikiye tümleyen & two's complement \\
ikilik taban & binary \\
%% isabet → bulma olarak B harfi altında (bkz. "bulma")
iş çıkarma oranı & throughput \\
iş parçacığı & thread \\
iş parçacığı düzeyi paralelliği & thread-level parallelism \\
iş parçacığı havuzu & thread pool \\
işaret biti & sign bit \\
işaretli genişletme & sign extension \\
işaretli & signed \\
işaretsiz & unsigned \\
işlem kodu & opcode \\
işlemci & processor \\
işletim sistemi & operating system \\
istisna (boru hattı) & exception (pipeline) \\

%% --- J ---
J türü & J-type \\

%% --- K ---
kritik yol & critical path \\
karşı bağımlılık & anti-dependence \\
karşılaştırma buyrukları & comparison instructions \\
kaydırma buyrukları & shift instructions \\
kayan nokta & floating point \\
kayan noktalı sayı birimi & floating point unit (\texttt{FPU}) \\
katı hal sürücü (\texttt{SSD}) & Solid State Drive (\texttt{SSD}) \\
kesin istisna & precise exception \\
kesme & interrupt \\
kesme denetleyicisi & interrupt controller \\
kesme hizmet yordamı & interrupt service routine \\
kilit & lock/mutex \\
kilitlenme & deadlock \\
kod çözücü & decoder \\
koşullu dallanma & conditional branch \\
koşulsuz dallanma & unconditional branch \\
kuantum hesaplama & quantum computing \\
kübit & qubit \\
küçüğü başta & little-endian \\
kırmık & die \\
küme ilişkili & set associative \\

%% --- L ---
L1 önbellek & L1 cache \\
L2 önbellek & L2 cache \\
L3 önbellek & L3 cache \\

%% --- M ---
\texttt{M.2} & \texttt{M.2} (form factor) \\
makine dili & machine language \\
makro buyruk & macro instruction \\
\texttt{MLC} & Multi-Level Cell \\
mantık buyrukları & logic instructions \\
mantık kapısı & logic gate \\
matris çarpımı & matrix multiplication \\
Meltdown & Meltdown \\
\texttt{MESI} protokolü & \texttt{MESI} protocol \\
merkezi işlem birimi (\texttt{MİB}) & Central Processing Unit (\texttt{CPU}) \\
mesaj geçirme & message passing \\
\texttt{MFLOPS} & Millions of Floating Point Operations Per Second \\
mikro buyruk & micro-operation \\
mikro program belleği & control store \\
mikromimari & microarchitecture \\
mikroprogramlama & microprogramming \\
\texttt{MIPS} & Millions of Instructions Per Second \\
\texttt{MİB} & \texttt{CPU} (Central Processing Unit) \\
Moore Yasası & Moore's Law \\

%% --- N ---
\texttt{NAND} Flaş bellek & \texttt{NAND Flash} memory \\
\texttt{NUMA} & Non-Uniform Memory Access \\
\texttt{NVMe} & Non-Volatile Memory Express \\

%% --- O ---
olağanlaştırma & normalization \\
on altılık taban & hexadecimal \\
optimal çıkarma algoritması & optimal (Bélády's) replacement algorithm \\
ortalama bellek erişim süresi (\texttt{AMAT}) & Average Memory Access Time (\texttt{AMAT}) \\

%% --- Ö ---
önbellek & cache \\
önbellek satırı & cache line \\
önceden getirme & prefetching \\
Özel VEYA kapısı & \texttt{XOR} gate \\
özel fonksiyon birimi (\texttt{SFU}) & Special Function Unit (\texttt{SFU}) \\

%% --- P ---
paralel hesaplama & parallel computing \\
paralel verimlilik & parallel efficiency \\
paylaşımlı bellek & shared memory \\
\texttt{PCI} Express & \texttt{PCI} Express \\
program sayacı & program counter \\

%% --- Q ---
\texttt{QLC} & Quad-Level Cell \\

%% --- R ---
R türü & R-type \\
\texttt{RAID} & Redundant Array of Independent Disks \\
\texttt{RISC} & Reduced Instruction Set Computer \\
\texttt{RISC-V} & \texttt{RISC-V} \\

%% --- S ---
S türü & S-type \\
sıfırla genişletme & zero extension \\
saat çevrimi & clock cycle \\
saat sıklığı & clock frequency \\
sabit bellek & constant memory \\
sabit disk sürücüsü (\texttt{HDD}) & hard disk drive (\texttt{HDD}) \\
\texttt{SATA} & Serial \texttt{ATA} \\
saklama buyruğu & store instruction \\
sanal adres & virtual address \\
sanal bellek & virtual memory \\
sayfa & page \\
sayfa hatası & page fault \\
sayfa tablosu & page table \\
sekizlik taban & octal \\
semafor & semaphore \\
sıra dışı yürütme & out-of-order execution \\
şeritleme & striping \\
\texttt{SIMD} & Single Instruction, Multiple Data \\
\texttt{SLC} & Single-Level Cell \\
\texttt{SIMT} & Single Instruction, Multiple Threads \\
Sinir Ağı İşlem Birimi (\texttt{NPU}) & Neural Processing Unit (\texttt{NPU}) \\
sınama programı (denektaşı) & benchmark \\
sistolik dizi & systolic array \\
Soğutma Tasarım Gücü (\texttt{TDP}) & Thermal Design Power (\texttt{TDP}) \\
\texttt{SoC} & System on Chip \\
sonlu durum makinesi & finite state machine \\
sonuç bağımlılığı & output dependence \\
sonradan yazma & write-back \\
soyutlama & abstraction \\
sözde buyruk & pseudo-instruction \\
sözde-EUZK & pseudo-LRU \\
Spectre & Spectre \\
\texttt{SPEC} & \texttt{SPEC} \\
spekülatif yürütme & speculative execution \\
\texttt{SRAM} & Static Random Access Memory \\
Sv39 (RISC-V) & Sv39 (39-bit virtual addressing) \\
Sv48 (RISC-V) & Sv48 (48-bit virtual addressing) \\
Sv57 (RISC-V) & Sv57 (57-bit virtual addressing) \\
çoklu buyruk başlatımlı & superscalar \\
süperpozisyon & superposition \\

%% --- Ş ---

%% --- T ---
tam ilişkili & fully associative \\
Thunderbolt & Thunderbolt \\
\texttt{TLC} & Triple-Level Cell \\
\texttt{TRIM} & \texttt{TRIM} \\
tam sayı & integer \\
tam toplayıcı & full adder \\
taşma & overflow \\
tek çevrimli & single-cycle \\
tek vuruşluk işlemci & single-cycle processor \\
tekdüze bellek erişimi & Uniform Memory Access (\texttt{UMA}) \\
tekdüze olmayan bellek erişimi & Non-Uniform Memory Access (\texttt{NUMA}) \\
Tensör çekirdeği & Tensor Core \\
Tensör İşlem Birimi (\texttt{TPU}) & Tensor Processing Unit (\texttt{TPU}) \\
\texttt{TLB} & Translation Lookaside Buffer \\
toplama & addition \\
transistör & transistor \\
turnuva öngörücüsü & tournament predictor \\
tümleşik devre & integrated circuit \\

%% --- U ---
uç bilgi işlem & edge computing \\
U türü & U-type \\
\texttt{USB} & Universal Serial Bus \\
\texttt{UCIe} & Universal Chiplet Interconnect Express \\
\texttt{UMA} & Uniform Memory Access \\
Unicode & Unicode \\
uzamsal yerellik & spatial locality \\

%% --- Ü ---
üç C modeli (3C modeli) & three Cs model \\
üs & exponent \\
üst anlık buyruklar & upper immediate instructions \\

%% --- V ---
vakum tüpü & vacuum tube \\
VE kapısı & \texttt{AND} gate \\
veri akış mimarisi & dataflow architecture \\
veri aktarma buyrukları & data transfer instructions \\
veri bağımlılığı & data dependency \\
veri belleği & data memory \\
veri tehlikesi & data hazard \\
veri yolu & datapath \\
veri yönlendirmesi & forwarding \\
verim & throughput \\
verim odaklı tasarım & throughput-oriented design \\
Verilog & Verilog \\
VEYA kapısı & \texttt{OR} gate \\
\texttt{VHDL} & \texttt{VHDL} \\
\texttt{VLSI} & Very Large-Scale Integration \\
Von Neumann darboğazı & Von Neumann bottleneck \\
Von Neumann mimarisi & Von Neumann architecture \\

%% --- Y ---
yan kanal saldırısı & side-channel attack \\
yanlış paylaşım & false sharing \\
yapay sınama programları & synthetic benchmarks \\
yapay zeka & artificial intelligence \\
yapısal tehlike & structural hazard \\
yarım toplayıcı & half adder \\
yarış durumu & race condition \\
yatay mikroprogramlama & horizontal microprogramming \\
yazılım & software \\
yazma geçersizleştirme & write-invalidate \\
yazma güncelleme & write-update \\
yazma serileştirmesi & write serialization \\
yazma tahsisli & write-allocate \\
yazma tahsissiz & no-write-allocate \\
yazma ara belleği & write buffer \\
yazma yayılımı & write propagation \\
yazmaç & register \\
yazmaç aktarım düzeyi (\texttt{RTL}) & Register Transfer Level (\texttt{RTL}) \\
yazmaç öbeği & register file \\
yazmaç yeniden adlandırma & register renaming \\
yekpare kırmık & monolithic die \\
yeniden yapılandırılabilir hesaplama & reconfigurable computing \\
yerel bellek & local memory \\
yerellik ilkesi & locality principle \\
yığıt & stack \\
yığıt göstergesi & stack pointer \\
yonga & chip \\
yuvarlama & rounding \\
yükleyici & loader \\
yükle-kullan sorunu & load-use hazard \\
yükleme buyruğu & load instruction \\
yürütme birimi & execution unit \\
yürütme zamanı & execution time \\

%% --- Z ---
zamansal yerellik & temporal locality \\

\end{longtable}
}


\bigskip

%% ============================================================
%%  BÖLÜM 2: İngilizce → Türkçe
%%  İngilizce alfabetik sıra (A-Z)
%% ============================================================
\section*{İngilizce -- Türkçe}
\addcontentsline{toc}{section}{İngilizce -- Türkçe}

{\small
\begin{longtable}{@{}p{7.6cm}p{7.6cm}@{}}
\toprule
\textbf{İngilizce Terim} & \textbf{Türkçe Karşılık} \\
\midrule
\endfirsthead

\toprule
\textbf{İngilizce Terim (cont.)} & \textbf{Türkçe Karşılık} \\
\midrule
\endhead

\midrule
\multicolumn{2}{r}{\small\textit{(continued on next page)}} \\
\endfoot

\bottomrule
\endlastfoot

%% --- A ---
abstraction & soyutlama \\
access time & erişim süresi \\
addition & toplama \\
Average Memory Access Time (\texttt{AMAT}) & ortalama bellek erişim süresi (\texttt{AMAT}) \\
addressing modes & adresleme kipleri \\
Amdahl's Law & Amdahl Yasası \\
\texttt{AND} gate & VE kapısı \\
anti-dependence & karşı bağımlılık \\
Application-Specific Integrated Circuit (\texttt{ASIC}) & \texttt{ASIC} \\
Arithmetic Logic Unit (\texttt{ALU}) & aritmetik mantık birimi (\texttt{AMB}) \\
arithmetic instructions & aritmetik buyruklar \\
arithmetic intensity & aritmetik yoğunluk \\
artificial intelligence & yapay zeka \\
\texttt{ASCII} & \texttt{ASCII} \\
assembler & çevirici \\
assembly language & çevirici dil \\
atomic operation & atomik işlem \\

%% --- B ---
B-type & B türü \\
Booth's algorithm & Booth algoritması \\
bandwidth & bant genişliği \\
barrier & engel \\
benchmark & sınama programı (denektaşı) \\
big-endian & büyüğü başta \\
big.LITTLE & big.LITTLE \\
binary & ikilik taban \\
Binary-Coded Decimal (\texttt{BCD}) & \texttt{BCD} \\
bit & bit \\
block & blok \\
bottleneck & darboğaz \\
branch instructions & dallanma buyrukları \\
branch prediction & dallanma öngörüsü \\
branch target buffer (BTB) & dallanma hedef ara belleği \\
buffer & ara bellek \\
bubble (pipeline) & boşluk (boru hattı) \\
bus & yol \\
byte & bayt \\
byte order & bayt sıralaması \\

%% --- C ---
cache & önbellek \\
cache coherence & bellek tutarlığı \\
cache line & önbellek satırı \\
calling conventions & çağrı kuralları \\
carry & elde \\
carry lookahead & elden önceden üretme \\
Central Processing Unit (\texttt{CPU}) & merkezi işlem birimi (\texttt{MİB}) \\
chip & yonga \\
chiplet & chiplet \\
clock cycle & saat çevrimi \\
clock frequency & saat sıklığı \\
combinational circuit & birleşik devre \\
Complementary \texttt{MOS} (\texttt{CMOS}) & \texttt{CMOS} \\
compiler & derleyici \\
comparison instructions & karşılaştırma buyrukları \\
Complex Instruction Set Computer (\texttt{CISC}) & \texttt{CISC} \\
Compute Unified Device Architecture (\texttt{CUDA}) & \texttt{CUDA} \\
concurrency & eşzamanlılık \\
conditional branch & koşullu dallanma \\
constant memory & sabit bellek \\
context switching & bağlam değiştirme \\
control hazard & denetim tehlikesi \\
control signal & denetim sinyali \\
control store & mikro program belleği \\
control table & denetim tablosu \\
control unit & denetim birimi \\
core & çekirdek \\
critical path & kritik yol \\
cycle & çevrim \\
cycle time & çevrim zamanı \\
cycles per instruction (\texttt{CPI}) & buyruk başına çevrim (\texttt{BBÇ}) \\

%% --- D ---
die & kırmık \\
data dependency & veri bağımlılığı \\
data hazard & veri tehlikesi \\
data memory & veri belleği \\
data transfer instructions & veri aktarma buyrukları \\
dataflow architecture & veri akış mimarisi \\
datapath & veri yolu \\
deadlock & kilitlenme \\
decoder & dekoder / kod çözücü \\
deep learning & derin öğrenme \\
delayed branching & geciktirilmiş dallanma \\
Dennard scaling & Dennard ölçeklemesi \\
device & aygıt \\
Direct Memory Access (\texttt{DMA}) & \texttt{DMA} \\
direct-mapped cache & doğrudan eşlemeli önbellek \\
directory-based protocol & dizin tabanlı protokol \\
distributed memory & dağıtık bellek \\
division & bölme \\
Double Data Rate (\texttt{DDR}) & \texttt{DDR} \\
dynamic prediction & devingen öngörü \\
dynamic power consumption & devingen güç tüketimi \\
Dynamic Random Access Memory (\texttt{DRAM}) & \texttt{DRAM} \\

%% --- E ---
edge computing & uç bilgi işlem \\
\texttt{EDVAC} & \texttt{EDVAC} \\
effective address & etkin adres \\
\texttt{ENIAC} & \texttt{ENIAC} \\
entanglement & dolanıklık \\
exception (pipeline) & istisna (boru hattı) \\
excess-N notation & artık-N gösterimi \\
execution time & yürütme zamanı \\
execution unit & yürütme birimi \\
exponent & üs \\

%% --- F ---
false sharing & yanlış paylaşım \\
fetch-decode-execute cycle & getir-çöz-yürüt döngüsü \\
Field-Programmable Gate Array (\texttt{FPGA}) & \texttt{FPGA} \\
finite state machine & sonlu durum makinesi \\
flash memory & flaş bellek \\
Flash Translation Layer (\texttt{FTL}) & Flaş çeviri katmanı (\texttt{FTL}) \\
flip-flop & flip-flop \\
floating point & kayan nokta \\
floating point unit (\texttt{FPU}) & kayan noktalı sayı birimi \\
forwarding & veri yönlendirmesi \\
frame pointer & çerçeve göstergesi \\
full adder & tam toplayıcı \\
fully associative & tam ilişkili \\

%% --- G ---
General-Purpose computing on GPU (\texttt{GPGPU}) & \texttt{GPGPU} \\
general-purpose register & genel amaçlı yazmaç \\
geometric mean & geometrik ortalama \\
global memory & genel bellek \\
Graphics Processing Unit (\texttt{GPU}) & \texttt{GPU} \\
grid & ızgara \\
gshare predictor & gshare öngörücü \\
Gustafson's Law & Gustafson Yasası \\

%% --- H ---
half adder & yarım toplayıcı \\
hard disk drive (\texttt{HDD}) & sabit disk sürücüsü \\
hardware & donanım \\
hardware multithreading & donanımsal çok kanallılık \\
Harvard architecture & Harvard mimarisi \\
hazard detection unit & düzenleme birimi \\
heterogeneous computing & heterojen hesaplama \\
hexadecimal & on altılık taban \\
High Bandwidth Memory (\texttt{HBM}) & \texttt{HBM} \\
hit (cache hit) & bulma (önbellekte) \\
hit rate & bulma oranı \\
horizontal microprogramming & yatay mikroprogramlama \\
huge page & büyük sayfa \\

%% --- I ---
I-type & I türü \\
\texttt{IEEE 754} & \texttt{IEEE 754} \\
immediate generator & anlık değer üreteci \\
immediate value & anlık değer \\
inference (AI) & çıkarım (yapay zeka) \\
Input/Output Operations Per Second (\texttt{IOPS}) & \texttt{IOPS} \\
input/output (\texttt{I/O}) & giriş/çıkış (\texttt{G/Ç}) \\
instruction & buyruk \\
instruction cache & buyruk önbelleği \\
instruction decode & buyruk çözme \\
instruction counter & buyruk sayacı \\
instruction fetch unit & buyruk getirme birimi \\
instruction format & buyruk biçimi \\
instruction memory & buyruk belleği \\
Instruction Set Architecture (\texttt{ISA}) & buyruk kümesi mimarisi \\
instruction-level parallelism (\texttt{ILP}) & buyruk düzeyi paralelliği \\
instructions per cycle (\texttt{IPC}) & çevrim başına buyruk (\texttt{ÇBB}) \\
integer & tam sayı \\
integrated circuit & tümleşik devre \\
interrupt & kesme \\
interrupt controller & kesme denetleyicisi \\
interrupt service routine & kesme hizmet yordamı \\

%% --- J ---
J-type & J türü \\
jump instructions & atlama buyrukları \\

%% --- K ---
kernel function (GPU) & çekirdek fonksiyonu (GPU) \\

%% --- L ---
L1 cache & L1 önbellek \\
L2 cache & L2 önbellek \\
L3 cache & L3 önbellek \\
latency & gecikme \\
latency-oriented design & gecikme odaklı tasarım \\
linker & bağlayıcı \\
little-endian & küçüğü başta \\
load instruction & yükleme buyruğu \\
load-use hazard & yükle-kullan sorunu \\
loader & yükleyici \\
local memory & yerel bellek \\
locality principle & yerellik ilkesi \\
Look-Up Table (\texttt{LUT}) & bakış tablosu (\texttt{LUT}) \\
lock/mutex & kilit \\
logic gate & mantık kapısı \\
logic instructions & mantık buyrukları \\

%% --- M ---
\texttt{M.2} & \texttt{M.2} (form faktörü) \\
machine language & makine dili \\
macro instruction & makro buyruk \\
matrix multiplication & matris çarpımı \\
Meltdown & Meltdown \\
memory & bellek \\
memory alignment & bellek hizalama \\
memory coalescing & bellek birleştirme \\
memory consistency model & bellek tutarlık modeli \\
memory controller & bellek denetleyicisi \\
memory hierarchy & bellek hiyerarşisi \\
memory protection & bellek koruması \\
memory wall & bellek duvarı \\
memory-mapped I/O & bellekle eşlemeli G/Ç \\
\texttt{MESI} protocol & \texttt{MESI} protokolü \\
message passing & mesaj geçirme \\
microarchitecture & mikromimari \\
micro-operation & mikro buyruk \\
microprogramming & mikroprogramlama \\
Millions of Floating Point Operations Per Second (\texttt{MFLOPS}) & \texttt{MFLOPS} \\
Millions of Instructions Per Second (\texttt{MIPS}) & \texttt{MIPS} \\
miss (cache miss) & bulamama (önbellekte) \\
miss penalty & bulamama cezası \\
miss rate & bulamama oranı \\
monolithic die & yekpare kırmık \\
Moore's Law & Moore Yasası \\
multi-core processor & çok çekirdekli işlemci \\
multiply-accumulate (\texttt{MAC}) & çarpma-toplama (\texttt{MAC}) \\
multi-cycle & çok çevrimli \\
multi-cycle processor & çok vuruşluk işlemci \\
Multi-Level Cell (\texttt{MLC}) & \texttt{MLC} \\
multi-level page table & çok seviyeli sayfa tablosu \\
multiplication & çarpma \\
multiplexer & çoklayıcı \\

%% --- N ---
\texttt{NAND Flash} memory & \texttt{NAND} Flaş bellek \\
Network on Chip (\texttt{NoC}) & yonga içi bağlantı ağı (\texttt{NoC}) \\
Neural Processing Unit (\texttt{NPU}) & Sinir Ağı İşlem Birimi (\texttt{NPU}) \\
Non-Uniform Memory Access (\texttt{NUMA}) & tekdüze olmayan bellek erişimi \\
Non-Volatile Memory Express (\texttt{NVMe}) & \texttt{NVMe} \\
no-write-allocate & yazma tahsissiz \\
normalization & olağanlaştırma \\
\texttt{NOT} gate & DEĞİL kapısı \\

%% --- O ---
occupancy & doluluk \\
octal & sekizlik taban \\
one's complement & bire tümleyen \\
opcode & işlem kodu \\
operating system & işletim sistemi \\
optimal (Bélády's) replacement algorithm & optimal çıkarma algoritması \\
\texttt{OR} gate & VEYA kapısı \\
out-of-order execution & sıra dışı yürütme \\
output dependence & sonuç bağımlılığı \\
overflow & taşma \\

%% --- P ---
page & sayfa \\
parity (RAID) & eşlik (RAID) \\
page fault & sayfa hatası \\
page table & sayfa tablosu \\
parallel computing & paralel hesaplama \\
parallel efficiency & paralel verimlilik \\
\texttt{PCI} Express & \texttt{PCI} Express \\
performance & başarım \\
pipeline & boru hattı \\
pipeline hazard & boru hattı tehlikesi \\
pipeline register & boru hattı kaydı \\
pipeline stage & boru hattı aşaması \\
port & bağlantı noktası \\
power consumption & güç tüketimi \\
power wall & güç duvarı \\
precise exception & kesin istisna \\
prefetching & önceden getirme \\
privileged mode & ayrıcalıklı kip \\
processor & işlemci \\
program counter & program sayacı \\
pseudo-instruction & sözde buyruk \\
pseudo-LRU & sözde-EUZK \\

%% --- Q ---
Quad-Level Cell (\texttt{QLC}) & \texttt{QLC} \\
quantum computing & kuantum hesaplama \\
qubit & kübit \\

%% --- R ---
R-type & R türü \\
race condition & yarış durumu \\
reconfigurable computing & yeniden yapılandırılabilir hesaplama \\
Redundant Array of Independent Disks (\texttt{RAID}) & \texttt{RAID} \\
Reduced Instruction Set Computer (\texttt{RISC}) & \texttt{RISC} \\
register & yazmaç \\
register file & yazmaç öbeği \\
register renaming & yazmaç yeniden adlandırma \\
Register Transfer Level (\texttt{RTL}) & yazmaç aktarım düzeyi (\texttt{RTL}) \\
relaxed consistency & gevşek tutarlılık \\
return address & dönüş adresi \\
ripple carry adder & dalga eldeli toplayıcı \\
\texttt{RISC-V} & \texttt{RISC-V} \\
roofline model & çatı çizgisi modeli \\
rotational latency & dönme gecikmesi \\
rounding & yuvarlama \\

%% --- S ---
S-type & S türü \\
\texttt{SATA} (Serial \texttt{ATA}) & \texttt{SATA} \\
secondary storage & ikincil bellek \\
seek time & arama zamanı \\
semaphore & semafor \\
sequential consistency & ardışık tutarlılık \\
Serial \texttt{ATA} (\texttt{SATA}) & \texttt{SATA} \\
set associative & küme ilişkili \\
shader & gölgelendirici \\
shared memory & paylaşımlı bellek \\
shift instructions & kaydırma buyrukları \\
side-channel attack & yan kanal saldırısı \\
sign bit & işaret biti \\
sign extension & işaretli genişletme \\
signed & işaretli \\
single-cycle processor & tek vuruşluk işlemci \\
Simultaneous Multithreading (\texttt{SMT}) & eşzamanlı çok kanallılık (\texttt{SMT}) \\
Single Instruction, Multiple Data (\texttt{SIMD}) & \texttt{SIMD} \\
Single Instruction, Multiple Threads (\texttt{SIMT}) & \texttt{SIMT} \\
Single-Level Cell (\texttt{SLC}) & \texttt{SLC} \\
single-cycle & tek çevrimli \\
snooping protocol & gözetleme protokolü \\
software & yazılım \\
scientific notation & bilimsel gösterim \\
significand / mantissa & anlamlı kısım \\
Solid State Drive (\texttt{SSD}) & katı hal sürücü (\texttt{SSD}) \\
spatial locality & uzamsal yerellik \\
\texttt{SPEC} & \texttt{SPEC} \\
Special Function Unit (\texttt{SFU}) & özel fonksiyon birimi (\texttt{SFU}) \\
Spectre & Spectre \\
speculative execution & spekülatif yürütme \\
spin-lock & döngüsel kilit \\
stack & yığıt \\
stack pointer & yığıt göstergesi \\
stale data & bayat veri \\
stall & durma \\
state register & durum yazmacı \\
static power consumption (leakage) & durağan güç tüketimi \\
static prediction & durağan öngörü \\
Static Random Access Memory (\texttt{SRAM}) & \texttt{SRAM} \\
store instruction & saklama buyruğu \\
Streaming Multiprocessor (\texttt{SM}) & Akış Çoklu İşlemci (\texttt{SM}) \\
Sv39 (39-bit virtual addressing) & Sv39 (RISC-V) \\
Sv48 (48-bit virtual addressing) & Sv48 (RISC-V) \\
Sv57 (57-bit virtual addressing) & Sv57 (RISC-V) \\
striping & şeritleme \\
structural hazard & yapısal tehlike \\
subroutine & altyordam \\
subtraction & çıkarma \\
superposition & süperpozisyon \\
superscalar & çoklu buyruk başlatımlı \\
synchronization & eşzamanlama \\
synthetic benchmarks & yapay sınama programları \\
System on Chip (\texttt{SoC}) & \texttt{SoC} \\
systolic array & sistolik dizi \\

%% --- T ---
tag & etiket \\
temporal locality & zamansal yerellik \\
three Cs model (3C) & üç C modeli (3C modeli) \\
Tensor Core & Tensör çekirdeği \\
Tensor Processing Unit (\texttt{TPU}) & Tensör İşlem Birimi (\texttt{TPU}) \\
texture memory & doku belleği \\
Thermal Design Power (\texttt{TDP}) & Soğutma Tasarım Gücü (\texttt{TDP}) \\
thread & iş parçacığı \\
thread pool & iş parçacığı havuzu \\
thread-level parallelism & iş parçacığı düzeyi paralelliği \\
throughput & verim / iş çıkarma oranı \\
Thunderbolt & Thunderbolt \\
training (AI) & eğitim (yapay zeka) \\
transistor & transistör \\
\texttt{TRIM} & \texttt{TRIM} \\
tournament predictor & turnuva öngörücüsü \\
Triple-Level Cell (\texttt{TLC}) & \texttt{TLC} \\
Translation Lookaside Buffer (\texttt{TLB}) & \texttt{TLB} \\
true data dependence & gerçek veri bağımlılığı \\
truth table & doğruluk tablosu \\
two's complement & ikiye tümleyen \\

%% --- U ---
U-type & U türü \\
unconditional branch & koşulsuz dallanma \\
Unicode & Unicode \\
Uniform Memory Access (\texttt{UMA}) & tekdüze bellek erişimi \\
unified memory architecture & birleşik bellek mimarisi \\
Universal Chiplet Interconnect Express (\texttt{UCIe}) & \texttt{UCIe} \\
Universal Serial Bus (\texttt{USB}) & \texttt{USB} \\
unsigned & işaretsiz \\
upper immediate instructions & üst anlık buyruklar \\

%% --- V ---
vacuum tube & vakum tüpü \\
Verilog & Verilog \\
vertical microprogramming & dikey mikroprogramlama \\
Very Large-Scale Integration (\texttt{VLSI}) & \texttt{VLSI} \\
\texttt{VHDL} & \texttt{VHDL} \\
virtual address & sanal adres \\
virtual memory & sanal bellek \\
Von Neumann architecture & Von Neumann mimarisi \\
Von Neumann bottleneck & Von Neumann darboğazı \\

%% --- W ---
warp & çözgü \\
weighted average & ağırlıklı ortalama \\
wear leveling & aşınma dengeleme \\
warp divergence & çözgü sapması \\
write-allocate & yazma tahsisli \\
write-back & geri yazma \\
write buffer & yazma ara belleği \\
write-invalidate & yazma geçersizleştirme \\
write propagation & yazma yayılımı \\
write serialization & yazma serileştirmesi \\
write-through & doğrudan yazma \\
write-update & yazma güncelleme \\

%% --- X ---
\texttt{XOR} gate & Özel VEYA kapısı \\

%% --- Z ---
zero extension & sıfırla genişletme \\

\end{longtable}
}
 ile dahil edilir.
% \chapter komutu ana-kitap.tex dosyasında tanımlanmalıdır.
% =============================================================

%% ============================================================
%%  BÖLÜM 1: Türkçe → İngilizce
%%  Türkçe alfabetik sıra: a,b,c,ç,d,e,f,g,ğ,h,ı,i,j,k,l,m,
%%                          n,o,ö,p,r,s,ş,t,u,ü,v,y,z
%% ============================================================
\section*{Türkçe -- İngilizce}
\addcontentsline{toc}{section}{Türkçe -- İngilizce}

{\small
\begin{longtable}{@{}p{7.6cm}p{7.6cm}@{}}
\toprule
\textbf{Türkçe Terim} & \textbf{İngilizce Karşılık} \\
\midrule
\endfirsthead

\toprule
\textbf{Türkçe Terim (devam)} & \textbf{İngilizce Karşılık} \\
\midrule
\endhead

\midrule
\multicolumn{2}{r}{\small\textit{(devamı sonraki sayfada)}} \\
\endfoot

\bottomrule
\endlastfoot

%% --- A ---
adresleme kipleri & addressing modes \\
ağırlıklı ortalama & weighted average \\
Akış Çoklu İşlemci (\texttt{SM}) & Streaming Multiprocessor (\texttt{SM}) \\
altyordam & subroutine \\
Amdahl Yasası & Amdahl's Law \\
\texttt{AMB} & \texttt{ALU} (Arithmetic Logic Unit) \\
anlamlı kısım & significand / mantissa \\
anlık değer & immediate value \\
anlık değer üreteci & immediate generator \\
ara bellek & buffer \\
ardışık tutarlılık & sequential consistency \\
arama zamanı & seek time \\
aritmetik buyruklar & arithmetic instructions \\
aşınma dengeleme & wear leveling \\
aritmetik mantık birimi (\texttt{AMB}) & Arithmetic Logic Unit (\texttt{ALU}) \\
aritmetik yoğunluk & arithmetic intensity \\
artık-N gösterimi & excess-N notation \\
\texttt{ASCII} & \texttt{ASCII} \\
\texttt{ASIC} & Application-Specific Integrated Circuit \\
çevirici dil (Assembly) & assembly language \\
atlama buyrukları & jump instructions \\
atomik işlem & atomic operation \\
aygıt & device \\
ayrıcalıklı kip & privileged mode \\

%% --- B ---
bilimsel gösterim & scientific notation \\
Booth algoritması & Booth's algorithm \\
B türü & B-type \\
bağlam değiştirme & context switching \\
bağlantı noktası & port \\
bağlayıcı & linker \\
bakış tablosu (\texttt{LUT}) & Look-Up Table (\texttt{LUT}) \\
bant genişliği & bandwidth \\
başarım & performance \\
bayat veri & stale data \\
bayt & byte \\
bayt sıralaması & byte order \\
büyük sayfa & huge page \\
\texttt{BBÇ} & \texttt{CPI} (cycles per instruction) \\
\texttt{BCD} & Binary-Coded Decimal \\
bellek & memory \\
bellek birleştirme & memory coalescing \\
bellek denetleyicisi & memory controller \\
bellek duvarı & memory wall \\
bellek koruması & memory protection \\
bellek hizalama & memory alignment \\
bellek hiyerarşisi & memory hierarchy \\
bellek tutarlığı & cache coherence \\
bellek tutarlık modeli & memory consistency model \\
bellekle eşlemeli G/Ç & memory-mapped I/O \\
benchmark & sınama programı (denektaşı) \\
big.LITTLE & big.LITTLE \\
bire tümleyen & one's complement \\
birleşik bellek mimarisi & unified memory architecture \\
birleşik devre & combinational circuit \\
bit & bit \\
blok & block \\
boru hattı & pipeline \\
boru hattı aşaması & pipeline stage \\
boru hattı kaydı & pipeline register \\
boru hattı tehlikesi & pipeline hazard \\
boşluk (boru hattı) & bubble / NOP \\
bölme & division \\
bulamama (önbellekte) & miss (cache miss) \\
bulamama cezası & miss penalty \\
bulamama oranı & miss rate \\
bulma (önbellekte) & hit (cache hit) \\
bulma oranı & hit rate \\
buyruk & instruction \\
buyruk başına çevrim (\texttt{BBÇ}) & cycles per instruction (\texttt{CPI}) \\
buyruk belleği & instruction memory \\
buyruk çözme & instruction decode \\
buyruk düzeyi paralelliği & instruction-level parallelism (\texttt{ILP}) \\
buyruk biçimi & instruction format \\
buyruk getirme birimi & instruction fetch unit \\
buyruk kümesi mimarisi & Instruction Set Architecture (\texttt{ISA}) \\
buyruk önbelleği & instruction cache \\
buyruk sayacı & instruction counter \\
büyüğü başta & big-endian \\

%% --- C ---
chiplet & chiplet \\
\texttt{CISC} & Complex Instruction Set Computer \\
\texttt{CMOS} & Complementary \texttt{MOS} \\
\texttt{CUDA} & Compute Unified Device Architecture \\

%% --- Ç ---
\texttt{ÇBB} & \texttt{IPC} (instructions per cycle) \\
çağrı kuralları & calling conventions \\
çarpma & multiplication \\
çarpma-toplama (\texttt{MAC}) & multiply-accumulate (\texttt{MAC}) \\
çatı çizgisi modeli & roofline model \\
çekirdek & core \\
çekirdek fonksiyonu (GPU) & kernel function \\
çerçeve göstergesi & frame pointer \\
çevirici & assembler \\
çevrim & cycle \\
çevrim başına buyruk (\texttt{ÇBB}) & instructions per cycle (\texttt{IPC}) \\
çevrim zamanı & cycle time \\
çıkarım (yapay zeka) & inference \\
çıkarma & subtraction \\
yonga içi bağlantı ağı (\texttt{NoC}) & Network on Chip (\texttt{NoC}) \\
çok çekirdekli işlemci & multi-core processor \\
çok seviyeli sayfa tablosu & multi-level page table \\
çok çevrimli & multi-cycle \\
çok vuruşluk işlemci & multi-cycle processor \\
çoklayıcı & multiplexer \\
çözgü & warp \\
çözgü sapması & warp divergence \\

%% --- D ---
dağıtık bellek & distributed memory \\
durum yazmacı & state register \\
dalga eldeli toplayıcı & ripple carry adder \\
dallanma buyrukları & branch instructions \\
dallanma hedef ara belleği & branch target buffer (BTB) \\
dallanma öngörüsü & branch prediction \\
darboğaz & bottleneck \\
\texttt{DDR} & Double Data Rate \\
DEĞİL kapısı & \texttt{NOT} gate \\
dekoder & decoder \\
denetim birimi & control unit \\
denetim sinyali & control signal \\
denetim tablosu & control table \\
denetim tehlikesi & control hazard \\
devingen güç tüketimi & dynamic power consumption \\
devingen öngörü & dynamic prediction \\
Dennard ölçeklemesi & Dennard scaling \\
derin öğrenme & deep learning \\
derleyici & compiler \\
dikey mikroprogramlama & vertical microprogramming \\
dizin tabanlı protokol & directory-based protocol \\
\texttt{DMA} & Direct Memory Access \\
\texttt{DRAM} & Dynamic Random Access Memory \\
doğrudan eşlemeli önbellek & direct-mapped cache \\
doğruluk tablosu & truth table \\
dönme gecikmesi & rotational latency \\
doğrudan yazma & write-through \\
doku belleği & texture memory \\
dolanıklık & entanglement \\
doluluk & occupancy \\
donanım & hardware \\
donanımsal çok kanallılık & hardware multithreading \\
dönüş adresi & return address \\
döngüsel kilit & spin-lock \\
durağan öngörü & static prediction \\
durağan güç tüketimi & static power consumption (leakage) \\
durma & stall \\
düzenleme birimi & hazard detection unit \\

%% --- E ---
\texttt{EDVAC} & \texttt{EDVAC} \\
eğitim (yapay zeka) & training \\
elde & carry \\
elden önceden üretme & carry lookahead \\
\texttt{ENIAC} & \texttt{ENIAC} \\
engel & barrier \\
eşlik (RAID) & parity (RAID) \\
erişim süresi & access time \\
eşzamanlama & synchronization \\
eşzamanlı çok kanallılık (\texttt{SMT}) & Simultaneous Multithreading (\texttt{SMT}) \\
eşzamanlılık & concurrency \\
etiket & tag \\
etkin adres & effective address \\

%% --- F ---
Flaş çeviri katmanı (\texttt{FTL}) & Flash Translation Layer (\texttt{FTL}) \\
flaş bellek & flash memory \\
flip-flop & flip-flop \\
\texttt{FPGA} & Field-Programmable Gate Array \\

%% --- G ---
gecikme & latency \\
geometrik ortalama & geometric mean \\
gecikme odaklı tasarım & latency-oriented design \\
geciktirilmiş dallanma & delayed branching \\
gerçek veri bağımlılığı & true data dependence \\
genel amaçlı yazmaç & general-purpose register \\
genel bellek & global memory \\
geri yazma & write-back \\
getir-çöz-yürüt döngüsü & fetch-decode-execute cycle \\
gevşek tutarlılık & relaxed consistency \\
giriş/çıkış (\texttt{G/Ç}) & input/output (\texttt{I/O}) \\
gölgelendirici & shader \\
gözetleme protokolü & snooping protocol \\
\texttt{GPGPU} & General-Purpose computing on GPU \\
\texttt{GPU} & Graphics Processing Unit \\
gshare öngörücü & gshare predictor \\
Gustafson Yasası & Gustafson's Law \\
güç duvarı & power wall \\
güç tüketimi & power consumption \\

%% --- Ğ ---

%% --- H ---
Harvard mimarisi & Harvard architecture \\
\texttt{HBM} & High Bandwidth Memory \\
heterojen hesaplama & heterogeneous computing \\

%% --- Önbellekte bulamama terimleri (B harfi altında da var) ---
%% --- I ---
ızgara & grid \\

%% --- İ ---
I türü & I-type \\
\texttt{IEEE 754} & \texttt{IEEE 754} \\
yol & bus \\
ikincil bellek & secondary storage \\
\texttt{IOPS} & Input/Output Operations Per Second \\
ikiye tümleyen & two's complement \\
ikilik taban & binary \\
%% isabet → bulma olarak B harfi altında (bkz. "bulma")
iş çıkarma oranı & throughput \\
iş parçacığı & thread \\
iş parçacığı düzeyi paralelliği & thread-level parallelism \\
iş parçacığı havuzu & thread pool \\
işaret biti & sign bit \\
işaretli genişletme & sign extension \\
işaretli & signed \\
işaretsiz & unsigned \\
işlem kodu & opcode \\
işlemci & processor \\
işletim sistemi & operating system \\
istisna (boru hattı) & exception (pipeline) \\

%% --- J ---
J türü & J-type \\

%% --- K ---
kritik yol & critical path \\
karşı bağımlılık & anti-dependence \\
karşılaştırma buyrukları & comparison instructions \\
kaydırma buyrukları & shift instructions \\
kayan nokta & floating point \\
kayan noktalı sayı birimi & floating point unit (\texttt{FPU}) \\
katı hal sürücü (\texttt{SSD}) & Solid State Drive (\texttt{SSD}) \\
kesin istisna & precise exception \\
kesme & interrupt \\
kesme denetleyicisi & interrupt controller \\
kesme hizmet yordamı & interrupt service routine \\
kilit & lock/mutex \\
kilitlenme & deadlock \\
kod çözücü & decoder \\
koşullu dallanma & conditional branch \\
koşulsuz dallanma & unconditional branch \\
kuantum hesaplama & quantum computing \\
kübit & qubit \\
küçüğü başta & little-endian \\
kırmık & die \\
küme ilişkili & set associative \\

%% --- L ---
L1 önbellek & L1 cache \\
L2 önbellek & L2 cache \\
L3 önbellek & L3 cache \\

%% --- M ---
\texttt{M.2} & \texttt{M.2} (form factor) \\
makine dili & machine language \\
makro buyruk & macro instruction \\
\texttt{MLC} & Multi-Level Cell \\
mantık buyrukları & logic instructions \\
mantık kapısı & logic gate \\
matris çarpımı & matrix multiplication \\
Meltdown & Meltdown \\
\texttt{MESI} protokolü & \texttt{MESI} protocol \\
merkezi işlem birimi (\texttt{MİB}) & Central Processing Unit (\texttt{CPU}) \\
mesaj geçirme & message passing \\
\texttt{MFLOPS} & Millions of Floating Point Operations Per Second \\
mikro buyruk & micro-operation \\
mikro program belleği & control store \\
mikromimari & microarchitecture \\
mikroprogramlama & microprogramming \\
\texttt{MIPS} & Millions of Instructions Per Second \\
\texttt{MİB} & \texttt{CPU} (Central Processing Unit) \\
Moore Yasası & Moore's Law \\

%% --- N ---
\texttt{NAND} Flaş bellek & \texttt{NAND Flash} memory \\
\texttt{NUMA} & Non-Uniform Memory Access \\
\texttt{NVMe} & Non-Volatile Memory Express \\

%% --- O ---
olağanlaştırma & normalization \\
on altılık taban & hexadecimal \\
optimal çıkarma algoritması & optimal (Bélády's) replacement algorithm \\
ortalama bellek erişim süresi (\texttt{AMAT}) & Average Memory Access Time (\texttt{AMAT}) \\

%% --- Ö ---
önbellek & cache \\
önbellek satırı & cache line \\
önceden getirme & prefetching \\
Özel VEYA kapısı & \texttt{XOR} gate \\
özel fonksiyon birimi (\texttt{SFU}) & Special Function Unit (\texttt{SFU}) \\

%% --- P ---
paralel hesaplama & parallel computing \\
paralel verimlilik & parallel efficiency \\
paylaşımlı bellek & shared memory \\
\texttt{PCI} Express & \texttt{PCI} Express \\
program sayacı & program counter \\

%% --- Q ---
\texttt{QLC} & Quad-Level Cell \\

%% --- R ---
R türü & R-type \\
\texttt{RAID} & Redundant Array of Independent Disks \\
\texttt{RISC} & Reduced Instruction Set Computer \\
\texttt{RISC-V} & \texttt{RISC-V} \\

%% --- S ---
S türü & S-type \\
sıfırla genişletme & zero extension \\
saat çevrimi & clock cycle \\
saat sıklığı & clock frequency \\
sabit bellek & constant memory \\
sabit disk sürücüsü (\texttt{HDD}) & hard disk drive (\texttt{HDD}) \\
\texttt{SATA} & Serial \texttt{ATA} \\
saklama buyruğu & store instruction \\
sanal adres & virtual address \\
sanal bellek & virtual memory \\
sayfa & page \\
sayfa hatası & page fault \\
sayfa tablosu & page table \\
sekizlik taban & octal \\
semafor & semaphore \\
sıra dışı yürütme & out-of-order execution \\
şeritleme & striping \\
\texttt{SIMD} & Single Instruction, Multiple Data \\
\texttt{SLC} & Single-Level Cell \\
\texttt{SIMT} & Single Instruction, Multiple Threads \\
Sinir Ağı İşlem Birimi (\texttt{NPU}) & Neural Processing Unit (\texttt{NPU}) \\
sınama programı (denektaşı) & benchmark \\
sistolik dizi & systolic array \\
Soğutma Tasarım Gücü (\texttt{TDP}) & Thermal Design Power (\texttt{TDP}) \\
\texttt{SoC} & System on Chip \\
sonlu durum makinesi & finite state machine \\
sonuç bağımlılığı & output dependence \\
sonradan yazma & write-back \\
soyutlama & abstraction \\
sözde buyruk & pseudo-instruction \\
sözde-EUZK & pseudo-LRU \\
Spectre & Spectre \\
\texttt{SPEC} & \texttt{SPEC} \\
spekülatif yürütme & speculative execution \\
\texttt{SRAM} & Static Random Access Memory \\
Sv39 (RISC-V) & Sv39 (39-bit virtual addressing) \\
Sv48 (RISC-V) & Sv48 (48-bit virtual addressing) \\
Sv57 (RISC-V) & Sv57 (57-bit virtual addressing) \\
çoklu buyruk başlatımlı & superscalar \\
süperpozisyon & superposition \\

%% --- Ş ---

%% --- T ---
tam ilişkili & fully associative \\
Thunderbolt & Thunderbolt \\
\texttt{TLC} & Triple-Level Cell \\
\texttt{TRIM} & \texttt{TRIM} \\
tam sayı & integer \\
tam toplayıcı & full adder \\
taşma & overflow \\
tek çevrimli & single-cycle \\
tek vuruşluk işlemci & single-cycle processor \\
tekdüze bellek erişimi & Uniform Memory Access (\texttt{UMA}) \\
tekdüze olmayan bellek erişimi & Non-Uniform Memory Access (\texttt{NUMA}) \\
Tensör çekirdeği & Tensor Core \\
Tensör İşlem Birimi (\texttt{TPU}) & Tensor Processing Unit (\texttt{TPU}) \\
\texttt{TLB} & Translation Lookaside Buffer \\
toplama & addition \\
transistör & transistor \\
turnuva öngörücüsü & tournament predictor \\
tümleşik devre & integrated circuit \\

%% --- U ---
uç bilgi işlem & edge computing \\
U türü & U-type \\
\texttt{USB} & Universal Serial Bus \\
\texttt{UCIe} & Universal Chiplet Interconnect Express \\
\texttt{UMA} & Uniform Memory Access \\
Unicode & Unicode \\
uzamsal yerellik & spatial locality \\

%% --- Ü ---
üç C modeli (3C modeli) & three Cs model \\
üs & exponent \\
üst anlık buyruklar & upper immediate instructions \\

%% --- V ---
vakum tüpü & vacuum tube \\
VE kapısı & \texttt{AND} gate \\
veri akış mimarisi & dataflow architecture \\
veri aktarma buyrukları & data transfer instructions \\
veri bağımlılığı & data dependency \\
veri belleği & data memory \\
veri tehlikesi & data hazard \\
veri yolu & datapath \\
veri yönlendirmesi & forwarding \\
verim & throughput \\
verim odaklı tasarım & throughput-oriented design \\
Verilog & Verilog \\
VEYA kapısı & \texttt{OR} gate \\
\texttt{VHDL} & \texttt{VHDL} \\
\texttt{VLSI} & Very Large-Scale Integration \\
Von Neumann darboğazı & Von Neumann bottleneck \\
Von Neumann mimarisi & Von Neumann architecture \\

%% --- Y ---
yan kanal saldırısı & side-channel attack \\
yanlış paylaşım & false sharing \\
yapay sınama programları & synthetic benchmarks \\
yapay zeka & artificial intelligence \\
yapısal tehlike & structural hazard \\
yarım toplayıcı & half adder \\
yarış durumu & race condition \\
yatay mikroprogramlama & horizontal microprogramming \\
yazılım & software \\
yazma geçersizleştirme & write-invalidate \\
yazma güncelleme & write-update \\
yazma serileştirmesi & write serialization \\
yazma tahsisli & write-allocate \\
yazma tahsissiz & no-write-allocate \\
yazma ara belleği & write buffer \\
yazma yayılımı & write propagation \\
yazmaç & register \\
yazmaç aktarım düzeyi (\texttt{RTL}) & Register Transfer Level (\texttt{RTL}) \\
yazmaç öbeği & register file \\
yazmaç yeniden adlandırma & register renaming \\
yekpare kırmık & monolithic die \\
yeniden yapılandırılabilir hesaplama & reconfigurable computing \\
yerel bellek & local memory \\
yerellik ilkesi & locality principle \\
yığıt & stack \\
yığıt göstergesi & stack pointer \\
yonga & chip \\
yuvarlama & rounding \\
yükleyici & loader \\
yükle-kullan sorunu & load-use hazard \\
yükleme buyruğu & load instruction \\
yürütme birimi & execution unit \\
yürütme zamanı & execution time \\

%% --- Z ---
zamansal yerellik & temporal locality \\

\end{longtable}
}


\bigskip

%% ============================================================
%%  BÖLÜM 2: İngilizce → Türkçe
%%  İngilizce alfabetik sıra (A-Z)
%% ============================================================
\section*{İngilizce -- Türkçe}
\addcontentsline{toc}{section}{İngilizce -- Türkçe}

{\small
\begin{longtable}{@{}p{7.6cm}p{7.6cm}@{}}
\toprule
\textbf{İngilizce Terim} & \textbf{Türkçe Karşılık} \\
\midrule
\endfirsthead

\toprule
\textbf{İngilizce Terim (cont.)} & \textbf{Türkçe Karşılık} \\
\midrule
\endhead

\midrule
\multicolumn{2}{r}{\small\textit{(continued on next page)}} \\
\endfoot

\bottomrule
\endlastfoot

%% --- A ---
abstraction & soyutlama \\
access time & erişim süresi \\
addition & toplama \\
Average Memory Access Time (\texttt{AMAT}) & ortalama bellek erişim süresi (\texttt{AMAT}) \\
addressing modes & adresleme kipleri \\
Amdahl's Law & Amdahl Yasası \\
\texttt{AND} gate & VE kapısı \\
anti-dependence & karşı bağımlılık \\
Application-Specific Integrated Circuit (\texttt{ASIC}) & \texttt{ASIC} \\
Arithmetic Logic Unit (\texttt{ALU}) & aritmetik mantık birimi (\texttt{AMB}) \\
arithmetic instructions & aritmetik buyruklar \\
arithmetic intensity & aritmetik yoğunluk \\
artificial intelligence & yapay zeka \\
\texttt{ASCII} & \texttt{ASCII} \\
assembler & çevirici \\
assembly language & çevirici dil \\
atomic operation & atomik işlem \\

%% --- B ---
B-type & B türü \\
Booth's algorithm & Booth algoritması \\
bandwidth & bant genişliği \\
barrier & engel \\
benchmark & sınama programı (denektaşı) \\
big-endian & büyüğü başta \\
big.LITTLE & big.LITTLE \\
binary & ikilik taban \\
Binary-Coded Decimal (\texttt{BCD}) & \texttt{BCD} \\
bit & bit \\
block & blok \\
bottleneck & darboğaz \\
branch instructions & dallanma buyrukları \\
branch prediction & dallanma öngörüsü \\
branch target buffer (BTB) & dallanma hedef ara belleği \\
buffer & ara bellek \\
bubble (pipeline) & boşluk (boru hattı) \\
bus & yol \\
byte & bayt \\
byte order & bayt sıralaması \\

%% --- C ---
cache & önbellek \\
cache coherence & bellek tutarlığı \\
cache line & önbellek satırı \\
calling conventions & çağrı kuralları \\
carry & elde \\
carry lookahead & elden önceden üretme \\
Central Processing Unit (\texttt{CPU}) & merkezi işlem birimi (\texttt{MİB}) \\
chip & yonga \\
chiplet & chiplet \\
clock cycle & saat çevrimi \\
clock frequency & saat sıklığı \\
combinational circuit & birleşik devre \\
Complementary \texttt{MOS} (\texttt{CMOS}) & \texttt{CMOS} \\
compiler & derleyici \\
comparison instructions & karşılaştırma buyrukları \\
Complex Instruction Set Computer (\texttt{CISC}) & \texttt{CISC} \\
Compute Unified Device Architecture (\texttt{CUDA}) & \texttt{CUDA} \\
concurrency & eşzamanlılık \\
conditional branch & koşullu dallanma \\
constant memory & sabit bellek \\
context switching & bağlam değiştirme \\
control hazard & denetim tehlikesi \\
control signal & denetim sinyali \\
control store & mikro program belleği \\
control table & denetim tablosu \\
control unit & denetim birimi \\
core & çekirdek \\
critical path & kritik yol \\
cycle & çevrim \\
cycle time & çevrim zamanı \\
cycles per instruction (\texttt{CPI}) & buyruk başına çevrim (\texttt{BBÇ}) \\

%% --- D ---
die & kırmık \\
data dependency & veri bağımlılığı \\
data hazard & veri tehlikesi \\
data memory & veri belleği \\
data transfer instructions & veri aktarma buyrukları \\
dataflow architecture & veri akış mimarisi \\
datapath & veri yolu \\
deadlock & kilitlenme \\
decoder & dekoder / kod çözücü \\
deep learning & derin öğrenme \\
delayed branching & geciktirilmiş dallanma \\
Dennard scaling & Dennard ölçeklemesi \\
device & aygıt \\
Direct Memory Access (\texttt{DMA}) & \texttt{DMA} \\
direct-mapped cache & doğrudan eşlemeli önbellek \\
directory-based protocol & dizin tabanlı protokol \\
distributed memory & dağıtık bellek \\
division & bölme \\
Double Data Rate (\texttt{DDR}) & \texttt{DDR} \\
dynamic prediction & devingen öngörü \\
dynamic power consumption & devingen güç tüketimi \\
Dynamic Random Access Memory (\texttt{DRAM}) & \texttt{DRAM} \\

%% --- E ---
edge computing & uç bilgi işlem \\
\texttt{EDVAC} & \texttt{EDVAC} \\
effective address & etkin adres \\
\texttt{ENIAC} & \texttt{ENIAC} \\
entanglement & dolanıklık \\
exception (pipeline) & istisna (boru hattı) \\
excess-N notation & artık-N gösterimi \\
execution time & yürütme zamanı \\
execution unit & yürütme birimi \\
exponent & üs \\

%% --- F ---
false sharing & yanlış paylaşım \\
fetch-decode-execute cycle & getir-çöz-yürüt döngüsü \\
Field-Programmable Gate Array (\texttt{FPGA}) & \texttt{FPGA} \\
finite state machine & sonlu durum makinesi \\
flash memory & flaş bellek \\
Flash Translation Layer (\texttt{FTL}) & Flaş çeviri katmanı (\texttt{FTL}) \\
flip-flop & flip-flop \\
floating point & kayan nokta \\
floating point unit (\texttt{FPU}) & kayan noktalı sayı birimi \\
forwarding & veri yönlendirmesi \\
frame pointer & çerçeve göstergesi \\
full adder & tam toplayıcı \\
fully associative & tam ilişkili \\

%% --- G ---
General-Purpose computing on GPU (\texttt{GPGPU}) & \texttt{GPGPU} \\
general-purpose register & genel amaçlı yazmaç \\
geometric mean & geometrik ortalama \\
global memory & genel bellek \\
Graphics Processing Unit (\texttt{GPU}) & \texttt{GPU} \\
grid & ızgara \\
gshare predictor & gshare öngörücü \\
Gustafson's Law & Gustafson Yasası \\

%% --- H ---
half adder & yarım toplayıcı \\
hard disk drive (\texttt{HDD}) & sabit disk sürücüsü \\
hardware & donanım \\
hardware multithreading & donanımsal çok kanallılık \\
Harvard architecture & Harvard mimarisi \\
hazard detection unit & düzenleme birimi \\
heterogeneous computing & heterojen hesaplama \\
hexadecimal & on altılık taban \\
High Bandwidth Memory (\texttt{HBM}) & \texttt{HBM} \\
hit (cache hit) & bulma (önbellekte) \\
hit rate & bulma oranı \\
horizontal microprogramming & yatay mikroprogramlama \\
huge page & büyük sayfa \\

%% --- I ---
I-type & I türü \\
\texttt{IEEE 754} & \texttt{IEEE 754} \\
immediate generator & anlık değer üreteci \\
immediate value & anlık değer \\
inference (AI) & çıkarım (yapay zeka) \\
Input/Output Operations Per Second (\texttt{IOPS}) & \texttt{IOPS} \\
input/output (\texttt{I/O}) & giriş/çıkış (\texttt{G/Ç}) \\
instruction & buyruk \\
instruction cache & buyruk önbelleği \\
instruction decode & buyruk çözme \\
instruction counter & buyruk sayacı \\
instruction fetch unit & buyruk getirme birimi \\
instruction format & buyruk biçimi \\
instruction memory & buyruk belleği \\
Instruction Set Architecture (\texttt{ISA}) & buyruk kümesi mimarisi \\
instruction-level parallelism (\texttt{ILP}) & buyruk düzeyi paralelliği \\
instructions per cycle (\texttt{IPC}) & çevrim başına buyruk (\texttt{ÇBB}) \\
integer & tam sayı \\
integrated circuit & tümleşik devre \\
interrupt & kesme \\
interrupt controller & kesme denetleyicisi \\
interrupt service routine & kesme hizmet yordamı \\

%% --- J ---
J-type & J türü \\
jump instructions & atlama buyrukları \\

%% --- K ---
kernel function (GPU) & çekirdek fonksiyonu (GPU) \\

%% --- L ---
L1 cache & L1 önbellek \\
L2 cache & L2 önbellek \\
L3 cache & L3 önbellek \\
latency & gecikme \\
latency-oriented design & gecikme odaklı tasarım \\
linker & bağlayıcı \\
little-endian & küçüğü başta \\
load instruction & yükleme buyruğu \\
load-use hazard & yükle-kullan sorunu \\
loader & yükleyici \\
local memory & yerel bellek \\
locality principle & yerellik ilkesi \\
Look-Up Table (\texttt{LUT}) & bakış tablosu (\texttt{LUT}) \\
lock/mutex & kilit \\
logic gate & mantık kapısı \\
logic instructions & mantık buyrukları \\

%% --- M ---
\texttt{M.2} & \texttt{M.2} (form faktörü) \\
machine language & makine dili \\
macro instruction & makro buyruk \\
matrix multiplication & matris çarpımı \\
Meltdown & Meltdown \\
memory & bellek \\
memory alignment & bellek hizalama \\
memory coalescing & bellek birleştirme \\
memory consistency model & bellek tutarlık modeli \\
memory controller & bellek denetleyicisi \\
memory hierarchy & bellek hiyerarşisi \\
memory protection & bellek koruması \\
memory wall & bellek duvarı \\
memory-mapped I/O & bellekle eşlemeli G/Ç \\
\texttt{MESI} protocol & \texttt{MESI} protokolü \\
message passing & mesaj geçirme \\
microarchitecture & mikromimari \\
micro-operation & mikro buyruk \\
microprogramming & mikroprogramlama \\
Millions of Floating Point Operations Per Second (\texttt{MFLOPS}) & \texttt{MFLOPS} \\
Millions of Instructions Per Second (\texttt{MIPS}) & \texttt{MIPS} \\
miss (cache miss) & bulamama (önbellekte) \\
miss penalty & bulamama cezası \\
miss rate & bulamama oranı \\
monolithic die & yekpare kırmık \\
Moore's Law & Moore Yasası \\
multi-core processor & çok çekirdekli işlemci \\
multiply-accumulate (\texttt{MAC}) & çarpma-toplama (\texttt{MAC}) \\
multi-cycle & çok çevrimli \\
multi-cycle processor & çok vuruşluk işlemci \\
Multi-Level Cell (\texttt{MLC}) & \texttt{MLC} \\
multi-level page table & çok seviyeli sayfa tablosu \\
multiplication & çarpma \\
multiplexer & çoklayıcı \\

%% --- N ---
\texttt{NAND Flash} memory & \texttt{NAND} Flaş bellek \\
Network on Chip (\texttt{NoC}) & yonga içi bağlantı ağı (\texttt{NoC}) \\
Neural Processing Unit (\texttt{NPU}) & Sinir Ağı İşlem Birimi (\texttt{NPU}) \\
Non-Uniform Memory Access (\texttt{NUMA}) & tekdüze olmayan bellek erişimi \\
Non-Volatile Memory Express (\texttt{NVMe}) & \texttt{NVMe} \\
no-write-allocate & yazma tahsissiz \\
normalization & olağanlaştırma \\
\texttt{NOT} gate & DEĞİL kapısı \\

%% --- O ---
occupancy & doluluk \\
octal & sekizlik taban \\
one's complement & bire tümleyen \\
opcode & işlem kodu \\
operating system & işletim sistemi \\
optimal (Bélády's) replacement algorithm & optimal çıkarma algoritması \\
\texttt{OR} gate & VEYA kapısı \\
out-of-order execution & sıra dışı yürütme \\
output dependence & sonuç bağımlılığı \\
overflow & taşma \\

%% --- P ---
page & sayfa \\
parity (RAID) & eşlik (RAID) \\
page fault & sayfa hatası \\
page table & sayfa tablosu \\
parallel computing & paralel hesaplama \\
parallel efficiency & paralel verimlilik \\
\texttt{PCI} Express & \texttt{PCI} Express \\
performance & başarım \\
pipeline & boru hattı \\
pipeline hazard & boru hattı tehlikesi \\
pipeline register & boru hattı kaydı \\
pipeline stage & boru hattı aşaması \\
port & bağlantı noktası \\
power consumption & güç tüketimi \\
power wall & güç duvarı \\
precise exception & kesin istisna \\
prefetching & önceden getirme \\
privileged mode & ayrıcalıklı kip \\
processor & işlemci \\
program counter & program sayacı \\
pseudo-instruction & sözde buyruk \\
pseudo-LRU & sözde-EUZK \\

%% --- Q ---
Quad-Level Cell (\texttt{QLC}) & \texttt{QLC} \\
quantum computing & kuantum hesaplama \\
qubit & kübit \\

%% --- R ---
R-type & R türü \\
race condition & yarış durumu \\
reconfigurable computing & yeniden yapılandırılabilir hesaplama \\
Redundant Array of Independent Disks (\texttt{RAID}) & \texttt{RAID} \\
Reduced Instruction Set Computer (\texttt{RISC}) & \texttt{RISC} \\
register & yazmaç \\
register file & yazmaç öbeği \\
register renaming & yazmaç yeniden adlandırma \\
Register Transfer Level (\texttt{RTL}) & yazmaç aktarım düzeyi (\texttt{RTL}) \\
relaxed consistency & gevşek tutarlılık \\
return address & dönüş adresi \\
ripple carry adder & dalga eldeli toplayıcı \\
\texttt{RISC-V} & \texttt{RISC-V} \\
roofline model & çatı çizgisi modeli \\
rotational latency & dönme gecikmesi \\
rounding & yuvarlama \\

%% --- S ---
S-type & S türü \\
\texttt{SATA} (Serial \texttt{ATA}) & \texttt{SATA} \\
secondary storage & ikincil bellek \\
seek time & arama zamanı \\
semaphore & semafor \\
sequential consistency & ardışık tutarlılık \\
Serial \texttt{ATA} (\texttt{SATA}) & \texttt{SATA} \\
set associative & küme ilişkili \\
shader & gölgelendirici \\
shared memory & paylaşımlı bellek \\
shift instructions & kaydırma buyrukları \\
side-channel attack & yan kanal saldırısı \\
sign bit & işaret biti \\
sign extension & işaretli genişletme \\
signed & işaretli \\
single-cycle processor & tek vuruşluk işlemci \\
Simultaneous Multithreading (\texttt{SMT}) & eşzamanlı çok kanallılık (\texttt{SMT}) \\
Single Instruction, Multiple Data (\texttt{SIMD}) & \texttt{SIMD} \\
Single Instruction, Multiple Threads (\texttt{SIMT}) & \texttt{SIMT} \\
Single-Level Cell (\texttt{SLC}) & \texttt{SLC} \\
single-cycle & tek çevrimli \\
snooping protocol & gözetleme protokolü \\
software & yazılım \\
scientific notation & bilimsel gösterim \\
significand / mantissa & anlamlı kısım \\
Solid State Drive (\texttt{SSD}) & katı hal sürücü (\texttt{SSD}) \\
spatial locality & uzamsal yerellik \\
\texttt{SPEC} & \texttt{SPEC} \\
Special Function Unit (\texttt{SFU}) & özel fonksiyon birimi (\texttt{SFU}) \\
Spectre & Spectre \\
speculative execution & spekülatif yürütme \\
spin-lock & döngüsel kilit \\
stack & yığıt \\
stack pointer & yığıt göstergesi \\
stale data & bayat veri \\
stall & durma \\
state register & durum yazmacı \\
static power consumption (leakage) & durağan güç tüketimi \\
static prediction & durağan öngörü \\
Static Random Access Memory (\texttt{SRAM}) & \texttt{SRAM} \\
store instruction & saklama buyruğu \\
Streaming Multiprocessor (\texttt{SM}) & Akış Çoklu İşlemci (\texttt{SM}) \\
Sv39 (39-bit virtual addressing) & Sv39 (RISC-V) \\
Sv48 (48-bit virtual addressing) & Sv48 (RISC-V) \\
Sv57 (57-bit virtual addressing) & Sv57 (RISC-V) \\
striping & şeritleme \\
structural hazard & yapısal tehlike \\
subroutine & altyordam \\
subtraction & çıkarma \\
superposition & süperpozisyon \\
superscalar & çoklu buyruk başlatımlı \\
synchronization & eşzamanlama \\
synthetic benchmarks & yapay sınama programları \\
System on Chip (\texttt{SoC}) & \texttt{SoC} \\
systolic array & sistolik dizi \\

%% --- T ---
tag & etiket \\
temporal locality & zamansal yerellik \\
three Cs model (3C) & üç C modeli (3C modeli) \\
Tensor Core & Tensör çekirdeği \\
Tensor Processing Unit (\texttt{TPU}) & Tensör İşlem Birimi (\texttt{TPU}) \\
texture memory & doku belleği \\
Thermal Design Power (\texttt{TDP}) & Soğutma Tasarım Gücü (\texttt{TDP}) \\
thread & iş parçacığı \\
thread pool & iş parçacığı havuzu \\
thread-level parallelism & iş parçacığı düzeyi paralelliği \\
throughput & verim / iş çıkarma oranı \\
Thunderbolt & Thunderbolt \\
training (AI) & eğitim (yapay zeka) \\
transistor & transistör \\
\texttt{TRIM} & \texttt{TRIM} \\
tournament predictor & turnuva öngörücüsü \\
Triple-Level Cell (\texttt{TLC}) & \texttt{TLC} \\
Translation Lookaside Buffer (\texttt{TLB}) & \texttt{TLB} \\
true data dependence & gerçek veri bağımlılığı \\
truth table & doğruluk tablosu \\
two's complement & ikiye tümleyen \\

%% --- U ---
U-type & U türü \\
unconditional branch & koşulsuz dallanma \\
Unicode & Unicode \\
Uniform Memory Access (\texttt{UMA}) & tekdüze bellek erişimi \\
unified memory architecture & birleşik bellek mimarisi \\
Universal Chiplet Interconnect Express (\texttt{UCIe}) & \texttt{UCIe} \\
Universal Serial Bus (\texttt{USB}) & \texttt{USB} \\
unsigned & işaretsiz \\
upper immediate instructions & üst anlık buyruklar \\

%% --- V ---
vacuum tube & vakum tüpü \\
Verilog & Verilog \\
vertical microprogramming & dikey mikroprogramlama \\
Very Large-Scale Integration (\texttt{VLSI}) & \texttt{VLSI} \\
\texttt{VHDL} & \texttt{VHDL} \\
virtual address & sanal adres \\
virtual memory & sanal bellek \\
Von Neumann architecture & Von Neumann mimarisi \\
Von Neumann bottleneck & Von Neumann darboğazı \\

%% --- W ---
warp & çözgü \\
weighted average & ağırlıklı ortalama \\
wear leveling & aşınma dengeleme \\
warp divergence & çözgü sapması \\
write-allocate & yazma tahsisli \\
write-back & geri yazma \\
write buffer & yazma ara belleği \\
write-invalidate & yazma geçersizleştirme \\
write propagation & yazma yayılımı \\
write serialization & yazma serileştirmesi \\
write-through & doğrudan yazma \\
write-update & yazma güncelleme \\

%% --- X ---
\texttt{XOR} gate & Özel VEYA kapısı \\

%% --- Z ---
zero extension & sıfırla genişletme \\

\end{longtable}
}
 ile dahil edilir.
% \chapter komutu ana-kitap.tex dosyasında tanımlanmalıdır.
% =============================================================

%% ============================================================
%%  BÖLÜM 1: Türkçe → İngilizce
%%  Türkçe alfabetik sıra: a,b,c,ç,d,e,f,g,ğ,h,ı,i,j,k,l,m,
%%                          n,o,ö,p,r,s,ş,t,u,ü,v,y,z
%% ============================================================
\section*{Türkçe -- İngilizce}
\addcontentsline{toc}{section}{Türkçe -- İngilizce}

{\small
\begin{longtable}{@{}p{7.6cm}p{7.6cm}@{}}
\toprule
\textbf{Türkçe Terim} & \textbf{İngilizce Karşılık} \\
\midrule
\endfirsthead

\toprule
\textbf{Türkçe Terim (devam)} & \textbf{İngilizce Karşılık} \\
\midrule
\endhead

\midrule
\multicolumn{2}{r}{\small\textit{(devamı sonraki sayfada)}} \\
\endfoot

\bottomrule
\endlastfoot

%% --- A ---
adresleme kipleri & addressing modes \\
ağırlıklı ortalama & weighted average \\
Akış Çoklu İşlemci (\texttt{SM}) & Streaming Multiprocessor (\texttt{SM}) \\
altyordam & subroutine \\
Amdahl Yasası & Amdahl's Law \\
\texttt{AMB} & \texttt{ALU} (Arithmetic Logic Unit) \\
anlamlı kısım & significand / mantissa \\
anlık değer & immediate value \\
anlık değer üreteci & immediate generator \\
ara bellek & buffer \\
ardışık tutarlılık & sequential consistency \\
arama zamanı & seek time \\
aritmetik buyruklar & arithmetic instructions \\
aşınma dengeleme & wear leveling \\
aritmetik mantık birimi (\texttt{AMB}) & Arithmetic Logic Unit (\texttt{ALU}) \\
aritmetik yoğunluk & arithmetic intensity \\
artık-N gösterimi & excess-N notation \\
\texttt{ASCII} & \texttt{ASCII} \\
\texttt{ASIC} & Application-Specific Integrated Circuit \\
çevirici dil (Assembly) & assembly language \\
atlama buyrukları & jump instructions \\
atomik işlem & atomic operation \\
aygıt & device \\
ayrıcalıklı kip & privileged mode \\

%% --- B ---
bilimsel gösterim & scientific notation \\
Booth algoritması & Booth's algorithm \\
B türü & B-type \\
bağlam değiştirme & context switching \\
bağlantı noktası & port \\
bağlayıcı & linker \\
bakış tablosu (\texttt{LUT}) & Look-Up Table (\texttt{LUT}) \\
bant genişliği & bandwidth \\
başarım & performance \\
bayat veri & stale data \\
bayt & byte \\
bayt sıralaması & byte order \\
büyük sayfa & huge page \\
\texttt{BBÇ} & \texttt{CPI} (cycles per instruction) \\
\texttt{BCD} & Binary-Coded Decimal \\
bellek & memory \\
bellek birleştirme & memory coalescing \\
bellek denetleyicisi & memory controller \\
bellek duvarı & memory wall \\
bellek koruması & memory protection \\
bellek hizalama & memory alignment \\
bellek hiyerarşisi & memory hierarchy \\
bellek tutarlığı & cache coherence \\
bellek tutarlık modeli & memory consistency model \\
bellekle eşlemeli G/Ç & memory-mapped I/O \\
benchmark & sınama programı (denektaşı) \\
big.LITTLE & big.LITTLE \\
bire tümleyen & one's complement \\
birleşik bellek mimarisi & unified memory architecture \\
birleşik devre & combinational circuit \\
bit & bit \\
blok & block \\
boru hattı & pipeline \\
boru hattı aşaması & pipeline stage \\
boru hattı kaydı & pipeline register \\
boru hattı tehlikesi & pipeline hazard \\
boşluk (boru hattı) & bubble / NOP \\
bölme & division \\
bulamama (önbellekte) & miss (cache miss) \\
bulamama cezası & miss penalty \\
bulamama oranı & miss rate \\
bulma (önbellekte) & hit (cache hit) \\
bulma oranı & hit rate \\
buyruk & instruction \\
buyruk başına çevrim (\texttt{BBÇ}) & cycles per instruction (\texttt{CPI}) \\
buyruk belleği & instruction memory \\
buyruk çözme & instruction decode \\
buyruk düzeyi paralelliği & instruction-level parallelism (\texttt{ILP}) \\
buyruk biçimi & instruction format \\
buyruk getirme birimi & instruction fetch unit \\
buyruk kümesi mimarisi & Instruction Set Architecture (\texttt{ISA}) \\
buyruk önbelleği & instruction cache \\
buyruk sayacı & instruction counter \\
büyüğü başta & big-endian \\

%% --- C ---
chiplet & chiplet \\
\texttt{CISC} & Complex Instruction Set Computer \\
\texttt{CMOS} & Complementary \texttt{MOS} \\
\texttt{CUDA} & Compute Unified Device Architecture \\

%% --- Ç ---
\texttt{ÇBB} & \texttt{IPC} (instructions per cycle) \\
çağrı kuralları & calling conventions \\
çarpma & multiplication \\
çarpma-toplama (\texttt{MAC}) & multiply-accumulate (\texttt{MAC}) \\
çatı çizgisi modeli & roofline model \\
çekirdek & core \\
çekirdek fonksiyonu (GPU) & kernel function \\
çerçeve göstergesi & frame pointer \\
çevirici & assembler \\
çevrim & cycle \\
çevrim başına buyruk (\texttt{ÇBB}) & instructions per cycle (\texttt{IPC}) \\
çevrim zamanı & cycle time \\
çıkarım (yapay zeka) & inference \\
çıkarma & subtraction \\
yonga içi bağlantı ağı (\texttt{NoC}) & Network on Chip (\texttt{NoC}) \\
çok çekirdekli işlemci & multi-core processor \\
çok seviyeli sayfa tablosu & multi-level page table \\
çok çevrimli & multi-cycle \\
çok vuruşluk işlemci & multi-cycle processor \\
çoklayıcı & multiplexer \\
çözgü & warp \\
çözgü sapması & warp divergence \\

%% --- D ---
dağıtık bellek & distributed memory \\
durum yazmacı & state register \\
dalga eldeli toplayıcı & ripple carry adder \\
dallanma buyrukları & branch instructions \\
dallanma hedef ara belleği & branch target buffer (BTB) \\
dallanma öngörüsü & branch prediction \\
darboğaz & bottleneck \\
\texttt{DDR} & Double Data Rate \\
DEĞİL kapısı & \texttt{NOT} gate \\
dekoder & decoder \\
denetim birimi & control unit \\
denetim sinyali & control signal \\
denetim tablosu & control table \\
denetim tehlikesi & control hazard \\
devingen güç tüketimi & dynamic power consumption \\
devingen öngörü & dynamic prediction \\
Dennard ölçeklemesi & Dennard scaling \\
derin öğrenme & deep learning \\
derleyici & compiler \\
dikey mikroprogramlama & vertical microprogramming \\
dizin tabanlı protokol & directory-based protocol \\
\texttt{DMA} & Direct Memory Access \\
\texttt{DRAM} & Dynamic Random Access Memory \\
doğrudan eşlemeli önbellek & direct-mapped cache \\
doğruluk tablosu & truth table \\
dönme gecikmesi & rotational latency \\
doğrudan yazma & write-through \\
doku belleği & texture memory \\
dolanıklık & entanglement \\
doluluk & occupancy \\
donanım & hardware \\
donanımsal çok kanallılık & hardware multithreading \\
dönüş adresi & return address \\
döngüsel kilit & spin-lock \\
durağan öngörü & static prediction \\
durağan güç tüketimi & static power consumption (leakage) \\
durma & stall \\
düzenleme birimi & hazard detection unit \\

%% --- E ---
\texttt{EDVAC} & \texttt{EDVAC} \\
eğitim (yapay zeka) & training \\
elde & carry \\
elden önceden üretme & carry lookahead \\
\texttt{ENIAC} & \texttt{ENIAC} \\
engel & barrier \\
eşlik (RAID) & parity (RAID) \\
erişim süresi & access time \\
eşzamanlama & synchronization \\
eşzamanlı çok kanallılık (\texttt{SMT}) & Simultaneous Multithreading (\texttt{SMT}) \\
eşzamanlılık & concurrency \\
etiket & tag \\
etkin adres & effective address \\

%% --- F ---
Flaş çeviri katmanı (\texttt{FTL}) & Flash Translation Layer (\texttt{FTL}) \\
flaş bellek & flash memory \\
flip-flop & flip-flop \\
\texttt{FPGA} & Field-Programmable Gate Array \\

%% --- G ---
gecikme & latency \\
geometrik ortalama & geometric mean \\
gecikme odaklı tasarım & latency-oriented design \\
geciktirilmiş dallanma & delayed branching \\
gerçek veri bağımlılığı & true data dependence \\
genel amaçlı yazmaç & general-purpose register \\
genel bellek & global memory \\
geri yazma & write-back \\
getir-çöz-yürüt döngüsü & fetch-decode-execute cycle \\
gevşek tutarlılık & relaxed consistency \\
giriş/çıkış (\texttt{G/Ç}) & input/output (\texttt{I/O}) \\
gölgelendirici & shader \\
gözetleme protokolü & snooping protocol \\
\texttt{GPGPU} & General-Purpose computing on GPU \\
\texttt{GPU} & Graphics Processing Unit \\
gshare öngörücü & gshare predictor \\
Gustafson Yasası & Gustafson's Law \\
güç duvarı & power wall \\
güç tüketimi & power consumption \\

%% --- Ğ ---

%% --- H ---
Harvard mimarisi & Harvard architecture \\
\texttt{HBM} & High Bandwidth Memory \\
heterojen hesaplama & heterogeneous computing \\

%% --- Önbellekte bulamama terimleri (B harfi altında da var) ---
%% --- I ---
ızgara & grid \\

%% --- İ ---
I türü & I-type \\
\texttt{IEEE 754} & \texttt{IEEE 754} \\
yol & bus \\
ikincil bellek & secondary storage \\
\texttt{IOPS} & Input/Output Operations Per Second \\
ikiye tümleyen & two's complement \\
ikilik taban & binary \\
%% isabet → bulma olarak B harfi altında (bkz. "bulma")
iş çıkarma oranı & throughput \\
iş parçacığı & thread \\
iş parçacığı düzeyi paralelliği & thread-level parallelism \\
iş parçacığı havuzu & thread pool \\
işaret biti & sign bit \\
işaretli genişletme & sign extension \\
işaretli & signed \\
işaretsiz & unsigned \\
işlem kodu & opcode \\
işlemci & processor \\
işletim sistemi & operating system \\
istisna (boru hattı) & exception (pipeline) \\

%% --- J ---
J türü & J-type \\

%% --- K ---
kritik yol & critical path \\
karşı bağımlılık & anti-dependence \\
karşılaştırma buyrukları & comparison instructions \\
kaydırma buyrukları & shift instructions \\
kayan nokta & floating point \\
kayan noktalı sayı birimi & floating point unit (\texttt{FPU}) \\
katı hal sürücü (\texttt{SSD}) & Solid State Drive (\texttt{SSD}) \\
kesin istisna & precise exception \\
kesme & interrupt \\
kesme denetleyicisi & interrupt controller \\
kesme hizmet yordamı & interrupt service routine \\
kilit & lock/mutex \\
kilitlenme & deadlock \\
kod çözücü & decoder \\
koşullu dallanma & conditional branch \\
koşulsuz dallanma & unconditional branch \\
kuantum hesaplama & quantum computing \\
kübit & qubit \\
küçüğü başta & little-endian \\
kırmık & die \\
küme ilişkili & set associative \\

%% --- L ---
L1 önbellek & L1 cache \\
L2 önbellek & L2 cache \\
L3 önbellek & L3 cache \\

%% --- M ---
\texttt{M.2} & \texttt{M.2} (form factor) \\
makine dili & machine language \\
makro buyruk & macro instruction \\
\texttt{MLC} & Multi-Level Cell \\
mantık buyrukları & logic instructions \\
mantık kapısı & logic gate \\
matris çarpımı & matrix multiplication \\
Meltdown & Meltdown \\
\texttt{MESI} protokolü & \texttt{MESI} protocol \\
merkezi işlem birimi (\texttt{MİB}) & Central Processing Unit (\texttt{CPU}) \\
mesaj geçirme & message passing \\
\texttt{MFLOPS} & Millions of Floating Point Operations Per Second \\
mikro buyruk & micro-operation \\
mikro program belleği & control store \\
mikromimari & microarchitecture \\
mikroprogramlama & microprogramming \\
\texttt{MIPS} & Millions of Instructions Per Second \\
\texttt{MİB} & \texttt{CPU} (Central Processing Unit) \\
Moore Yasası & Moore's Law \\

%% --- N ---
\texttt{NAND} Flaş bellek & \texttt{NAND Flash} memory \\
\texttt{NUMA} & Non-Uniform Memory Access \\
\texttt{NVMe} & Non-Volatile Memory Express \\

%% --- O ---
olağanlaştırma & normalization \\
on altılık taban & hexadecimal \\
optimal çıkarma algoritması & optimal (Bélády's) replacement algorithm \\
ortalama bellek erişim süresi (\texttt{AMAT}) & Average Memory Access Time (\texttt{AMAT}) \\

%% --- Ö ---
önbellek & cache \\
önbellek satırı & cache line \\
önceden getirme & prefetching \\
Özel VEYA kapısı & \texttt{XOR} gate \\
özel fonksiyon birimi (\texttt{SFU}) & Special Function Unit (\texttt{SFU}) \\

%% --- P ---
paralel hesaplama & parallel computing \\
paralel verimlilik & parallel efficiency \\
paylaşımlı bellek & shared memory \\
\texttt{PCI} Express & \texttt{PCI} Express \\
program sayacı & program counter \\

%% --- Q ---
\texttt{QLC} & Quad-Level Cell \\

%% --- R ---
R türü & R-type \\
\texttt{RAID} & Redundant Array of Independent Disks \\
\texttt{RISC} & Reduced Instruction Set Computer \\
\texttt{RISC-V} & \texttt{RISC-V} \\

%% --- S ---
S türü & S-type \\
sıfırla genişletme & zero extension \\
saat çevrimi & clock cycle \\
saat sıklığı & clock frequency \\
sabit bellek & constant memory \\
sabit disk sürücüsü (\texttt{HDD}) & hard disk drive (\texttt{HDD}) \\
\texttt{SATA} & Serial \texttt{ATA} \\
saklama buyruğu & store instruction \\
sanal adres & virtual address \\
sanal bellek & virtual memory \\
sayfa & page \\
sayfa hatası & page fault \\
sayfa tablosu & page table \\
sekizlik taban & octal \\
semafor & semaphore \\
sıra dışı yürütme & out-of-order execution \\
şeritleme & striping \\
\texttt{SIMD} & Single Instruction, Multiple Data \\
\texttt{SLC} & Single-Level Cell \\
\texttt{SIMT} & Single Instruction, Multiple Threads \\
Sinir Ağı İşlem Birimi (\texttt{NPU}) & Neural Processing Unit (\texttt{NPU}) \\
sınama programı (denektaşı) & benchmark \\
sistolik dizi & systolic array \\
Soğutma Tasarım Gücü (\texttt{TDP}) & Thermal Design Power (\texttt{TDP}) \\
\texttt{SoC} & System on Chip \\
sonlu durum makinesi & finite state machine \\
sonuç bağımlılığı & output dependence \\
sonradan yazma & write-back \\
soyutlama & abstraction \\
sözde buyruk & pseudo-instruction \\
sözde-EUZK & pseudo-LRU \\
Spectre & Spectre \\
\texttt{SPEC} & \texttt{SPEC} \\
spekülatif yürütme & speculative execution \\
\texttt{SRAM} & Static Random Access Memory \\
Sv39 (RISC-V) & Sv39 (39-bit virtual addressing) \\
Sv48 (RISC-V) & Sv48 (48-bit virtual addressing) \\
Sv57 (RISC-V) & Sv57 (57-bit virtual addressing) \\
çoklu buyruk başlatımlı & superscalar \\
süperpozisyon & superposition \\

%% --- Ş ---

%% --- T ---
tam ilişkili & fully associative \\
Thunderbolt & Thunderbolt \\
\texttt{TLC} & Triple-Level Cell \\
\texttt{TRIM} & \texttt{TRIM} \\
tam sayı & integer \\
tam toplayıcı & full adder \\
taşma & overflow \\
tek çevrimli & single-cycle \\
tek vuruşluk işlemci & single-cycle processor \\
tekdüze bellek erişimi & Uniform Memory Access (\texttt{UMA}) \\
tekdüze olmayan bellek erişimi & Non-Uniform Memory Access (\texttt{NUMA}) \\
Tensör çekirdeği & Tensor Core \\
Tensör İşlem Birimi (\texttt{TPU}) & Tensor Processing Unit (\texttt{TPU}) \\
\texttt{TLB} & Translation Lookaside Buffer \\
toplama & addition \\
transistör & transistor \\
turnuva öngörücüsü & tournament predictor \\
tümleşik devre & integrated circuit \\

%% --- U ---
uç bilgi işlem & edge computing \\
U türü & U-type \\
\texttt{USB} & Universal Serial Bus \\
\texttt{UCIe} & Universal Chiplet Interconnect Express \\
\texttt{UMA} & Uniform Memory Access \\
Unicode & Unicode \\
uzamsal yerellik & spatial locality \\

%% --- Ü ---
üç C modeli (3C modeli) & three Cs model \\
üs & exponent \\
üst anlık buyruklar & upper immediate instructions \\

%% --- V ---
vakum tüpü & vacuum tube \\
VE kapısı & \texttt{AND} gate \\
veri akış mimarisi & dataflow architecture \\
veri aktarma buyrukları & data transfer instructions \\
veri bağımlılığı & data dependency \\
veri belleği & data memory \\
veri tehlikesi & data hazard \\
veri yolu & datapath \\
veri yönlendirmesi & forwarding \\
verim & throughput \\
verim odaklı tasarım & throughput-oriented design \\
Verilog & Verilog \\
VEYA kapısı & \texttt{OR} gate \\
\texttt{VHDL} & \texttt{VHDL} \\
\texttt{VLSI} & Very Large-Scale Integration \\
Von Neumann darboğazı & Von Neumann bottleneck \\
Von Neumann mimarisi & Von Neumann architecture \\

%% --- Y ---
yan kanal saldırısı & side-channel attack \\
yanlış paylaşım & false sharing \\
yapay sınama programları & synthetic benchmarks \\
yapay zeka & artificial intelligence \\
yapısal tehlike & structural hazard \\
yarım toplayıcı & half adder \\
yarış durumu & race condition \\
yatay mikroprogramlama & horizontal microprogramming \\
yazılım & software \\
yazma geçersizleştirme & write-invalidate \\
yazma güncelleme & write-update \\
yazma serileştirmesi & write serialization \\
yazma tahsisli & write-allocate \\
yazma tahsissiz & no-write-allocate \\
yazma ara belleği & write buffer \\
yazma yayılımı & write propagation \\
yazmaç & register \\
yazmaç aktarım düzeyi (\texttt{RTL}) & Register Transfer Level (\texttt{RTL}) \\
yazmaç öbeği & register file \\
yazmaç yeniden adlandırma & register renaming \\
yekpare kırmık & monolithic die \\
yeniden yapılandırılabilir hesaplama & reconfigurable computing \\
yerel bellek & local memory \\
yerellik ilkesi & locality principle \\
yığıt & stack \\
yığıt göstergesi & stack pointer \\
yonga & chip \\
yuvarlama & rounding \\
yükleyici & loader \\
yükle-kullan sorunu & load-use hazard \\
yükleme buyruğu & load instruction \\
yürütme birimi & execution unit \\
yürütme zamanı & execution time \\

%% --- Z ---
zamansal yerellik & temporal locality \\

\end{longtable}
}


\bigskip

%% ============================================================
%%  BÖLÜM 2: İngilizce → Türkçe
%%  İngilizce alfabetik sıra (A-Z)
%% ============================================================
\section*{İngilizce -- Türkçe}
\addcontentsline{toc}{section}{İngilizce -- Türkçe}

{\small
\begin{longtable}{@{}p{7.6cm}p{7.6cm}@{}}
\toprule
\textbf{İngilizce Terim} & \textbf{Türkçe Karşılık} \\
\midrule
\endfirsthead

\toprule
\textbf{İngilizce Terim (cont.)} & \textbf{Türkçe Karşılık} \\
\midrule
\endhead

\midrule
\multicolumn{2}{r}{\small\textit{(continued on next page)}} \\
\endfoot

\bottomrule
\endlastfoot

%% --- A ---
abstraction & soyutlama \\
access time & erişim süresi \\
addition & toplama \\
Average Memory Access Time (\texttt{AMAT}) & ortalama bellek erişim süresi (\texttt{AMAT}) \\
addressing modes & adresleme kipleri \\
Amdahl's Law & Amdahl Yasası \\
\texttt{AND} gate & VE kapısı \\
anti-dependence & karşı bağımlılık \\
Application-Specific Integrated Circuit (\texttt{ASIC}) & \texttt{ASIC} \\
Arithmetic Logic Unit (\texttt{ALU}) & aritmetik mantık birimi (\texttt{AMB}) \\
arithmetic instructions & aritmetik buyruklar \\
arithmetic intensity & aritmetik yoğunluk \\
artificial intelligence & yapay zeka \\
\texttt{ASCII} & \texttt{ASCII} \\
assembler & çevirici \\
assembly language & çevirici dil \\
atomic operation & atomik işlem \\

%% --- B ---
B-type & B türü \\
Booth's algorithm & Booth algoritması \\
bandwidth & bant genişliği \\
barrier & engel \\
benchmark & sınama programı (denektaşı) \\
big-endian & büyüğü başta \\
big.LITTLE & big.LITTLE \\
binary & ikilik taban \\
Binary-Coded Decimal (\texttt{BCD}) & \texttt{BCD} \\
bit & bit \\
block & blok \\
bottleneck & darboğaz \\
branch instructions & dallanma buyrukları \\
branch prediction & dallanma öngörüsü \\
branch target buffer (BTB) & dallanma hedef ara belleği \\
buffer & ara bellek \\
bubble (pipeline) & boşluk (boru hattı) \\
bus & yol \\
byte & bayt \\
byte order & bayt sıralaması \\

%% --- C ---
cache & önbellek \\
cache coherence & bellek tutarlığı \\
cache line & önbellek satırı \\
calling conventions & çağrı kuralları \\
carry & elde \\
carry lookahead & elden önceden üretme \\
Central Processing Unit (\texttt{CPU}) & merkezi işlem birimi (\texttt{MİB}) \\
chip & yonga \\
chiplet & chiplet \\
clock cycle & saat çevrimi \\
clock frequency & saat sıklığı \\
combinational circuit & birleşik devre \\
Complementary \texttt{MOS} (\texttt{CMOS}) & \texttt{CMOS} \\
compiler & derleyici \\
comparison instructions & karşılaştırma buyrukları \\
Complex Instruction Set Computer (\texttt{CISC}) & \texttt{CISC} \\
Compute Unified Device Architecture (\texttt{CUDA}) & \texttt{CUDA} \\
concurrency & eşzamanlılık \\
conditional branch & koşullu dallanma \\
constant memory & sabit bellek \\
context switching & bağlam değiştirme \\
control hazard & denetim tehlikesi \\
control signal & denetim sinyali \\
control store & mikro program belleği \\
control table & denetim tablosu \\
control unit & denetim birimi \\
core & çekirdek \\
critical path & kritik yol \\
cycle & çevrim \\
cycle time & çevrim zamanı \\
cycles per instruction (\texttt{CPI}) & buyruk başına çevrim (\texttt{BBÇ}) \\

%% --- D ---
die & kırmık \\
data dependency & veri bağımlılığı \\
data hazard & veri tehlikesi \\
data memory & veri belleği \\
data transfer instructions & veri aktarma buyrukları \\
dataflow architecture & veri akış mimarisi \\
datapath & veri yolu \\
deadlock & kilitlenme \\
decoder & dekoder / kod çözücü \\
deep learning & derin öğrenme \\
delayed branching & geciktirilmiş dallanma \\
Dennard scaling & Dennard ölçeklemesi \\
device & aygıt \\
Direct Memory Access (\texttt{DMA}) & \texttt{DMA} \\
direct-mapped cache & doğrudan eşlemeli önbellek \\
directory-based protocol & dizin tabanlı protokol \\
distributed memory & dağıtık bellek \\
division & bölme \\
Double Data Rate (\texttt{DDR}) & \texttt{DDR} \\
dynamic prediction & devingen öngörü \\
dynamic power consumption & devingen güç tüketimi \\
Dynamic Random Access Memory (\texttt{DRAM}) & \texttt{DRAM} \\

%% --- E ---
edge computing & uç bilgi işlem \\
\texttt{EDVAC} & \texttt{EDVAC} \\
effective address & etkin adres \\
\texttt{ENIAC} & \texttt{ENIAC} \\
entanglement & dolanıklık \\
exception (pipeline) & istisna (boru hattı) \\
excess-N notation & artık-N gösterimi \\
execution time & yürütme zamanı \\
execution unit & yürütme birimi \\
exponent & üs \\

%% --- F ---
false sharing & yanlış paylaşım \\
fetch-decode-execute cycle & getir-çöz-yürüt döngüsü \\
Field-Programmable Gate Array (\texttt{FPGA}) & \texttt{FPGA} \\
finite state machine & sonlu durum makinesi \\
flash memory & flaş bellek \\
Flash Translation Layer (\texttt{FTL}) & Flaş çeviri katmanı (\texttt{FTL}) \\
flip-flop & flip-flop \\
floating point & kayan nokta \\
floating point unit (\texttt{FPU}) & kayan noktalı sayı birimi \\
forwarding & veri yönlendirmesi \\
frame pointer & çerçeve göstergesi \\
full adder & tam toplayıcı \\
fully associative & tam ilişkili \\

%% --- G ---
General-Purpose computing on GPU (\texttt{GPGPU}) & \texttt{GPGPU} \\
general-purpose register & genel amaçlı yazmaç \\
geometric mean & geometrik ortalama \\
global memory & genel bellek \\
Graphics Processing Unit (\texttt{GPU}) & \texttt{GPU} \\
grid & ızgara \\
gshare predictor & gshare öngörücü \\
Gustafson's Law & Gustafson Yasası \\

%% --- H ---
half adder & yarım toplayıcı \\
hard disk drive (\texttt{HDD}) & sabit disk sürücüsü \\
hardware & donanım \\
hardware multithreading & donanımsal çok kanallılık \\
Harvard architecture & Harvard mimarisi \\
hazard detection unit & düzenleme birimi \\
heterogeneous computing & heterojen hesaplama \\
hexadecimal & on altılık taban \\
High Bandwidth Memory (\texttt{HBM}) & \texttt{HBM} \\
hit (cache hit) & bulma (önbellekte) \\
hit rate & bulma oranı \\
horizontal microprogramming & yatay mikroprogramlama \\
huge page & büyük sayfa \\

%% --- I ---
I-type & I türü \\
\texttt{IEEE 754} & \texttt{IEEE 754} \\
immediate generator & anlık değer üreteci \\
immediate value & anlık değer \\
inference (AI) & çıkarım (yapay zeka) \\
Input/Output Operations Per Second (\texttt{IOPS}) & \texttt{IOPS} \\
input/output (\texttt{I/O}) & giriş/çıkış (\texttt{G/Ç}) \\
instruction & buyruk \\
instruction cache & buyruk önbelleği \\
instruction decode & buyruk çözme \\
instruction counter & buyruk sayacı \\
instruction fetch unit & buyruk getirme birimi \\
instruction format & buyruk biçimi \\
instruction memory & buyruk belleği \\
Instruction Set Architecture (\texttt{ISA}) & buyruk kümesi mimarisi \\
instruction-level parallelism (\texttt{ILP}) & buyruk düzeyi paralelliği \\
instructions per cycle (\texttt{IPC}) & çevrim başına buyruk (\texttt{ÇBB}) \\
integer & tam sayı \\
integrated circuit & tümleşik devre \\
interrupt & kesme \\
interrupt controller & kesme denetleyicisi \\
interrupt service routine & kesme hizmet yordamı \\

%% --- J ---
J-type & J türü \\
jump instructions & atlama buyrukları \\

%% --- K ---
kernel function (GPU) & çekirdek fonksiyonu (GPU) \\

%% --- L ---
L1 cache & L1 önbellek \\
L2 cache & L2 önbellek \\
L3 cache & L3 önbellek \\
latency & gecikme \\
latency-oriented design & gecikme odaklı tasarım \\
linker & bağlayıcı \\
little-endian & küçüğü başta \\
load instruction & yükleme buyruğu \\
load-use hazard & yükle-kullan sorunu \\
loader & yükleyici \\
local memory & yerel bellek \\
locality principle & yerellik ilkesi \\
Look-Up Table (\texttt{LUT}) & bakış tablosu (\texttt{LUT}) \\
lock/mutex & kilit \\
logic gate & mantık kapısı \\
logic instructions & mantık buyrukları \\

%% --- M ---
\texttt{M.2} & \texttt{M.2} (form faktörü) \\
machine language & makine dili \\
macro instruction & makro buyruk \\
matrix multiplication & matris çarpımı \\
Meltdown & Meltdown \\
memory & bellek \\
memory alignment & bellek hizalama \\
memory coalescing & bellek birleştirme \\
memory consistency model & bellek tutarlık modeli \\
memory controller & bellek denetleyicisi \\
memory hierarchy & bellek hiyerarşisi \\
memory protection & bellek koruması \\
memory wall & bellek duvarı \\
memory-mapped I/O & bellekle eşlemeli G/Ç \\
\texttt{MESI} protocol & \texttt{MESI} protokolü \\
message passing & mesaj geçirme \\
microarchitecture & mikromimari \\
micro-operation & mikro buyruk \\
microprogramming & mikroprogramlama \\
Millions of Floating Point Operations Per Second (\texttt{MFLOPS}) & \texttt{MFLOPS} \\
Millions of Instructions Per Second (\texttt{MIPS}) & \texttt{MIPS} \\
miss (cache miss) & bulamama (önbellekte) \\
miss penalty & bulamama cezası \\
miss rate & bulamama oranı \\
monolithic die & yekpare kırmık \\
Moore's Law & Moore Yasası \\
multi-core processor & çok çekirdekli işlemci \\
multiply-accumulate (\texttt{MAC}) & çarpma-toplama (\texttt{MAC}) \\
multi-cycle & çok çevrimli \\
multi-cycle processor & çok vuruşluk işlemci \\
Multi-Level Cell (\texttt{MLC}) & \texttt{MLC} \\
multi-level page table & çok seviyeli sayfa tablosu \\
multiplication & çarpma \\
multiplexer & çoklayıcı \\

%% --- N ---
\texttt{NAND Flash} memory & \texttt{NAND} Flaş bellek \\
Network on Chip (\texttt{NoC}) & yonga içi bağlantı ağı (\texttt{NoC}) \\
Neural Processing Unit (\texttt{NPU}) & Sinir Ağı İşlem Birimi (\texttt{NPU}) \\
Non-Uniform Memory Access (\texttt{NUMA}) & tekdüze olmayan bellek erişimi \\
Non-Volatile Memory Express (\texttt{NVMe}) & \texttt{NVMe} \\
no-write-allocate & yazma tahsissiz \\
normalization & olağanlaştırma \\
\texttt{NOT} gate & DEĞİL kapısı \\

%% --- O ---
occupancy & doluluk \\
octal & sekizlik taban \\
one's complement & bire tümleyen \\
opcode & işlem kodu \\
operating system & işletim sistemi \\
optimal (Bélády's) replacement algorithm & optimal çıkarma algoritması \\
\texttt{OR} gate & VEYA kapısı \\
out-of-order execution & sıra dışı yürütme \\
output dependence & sonuç bağımlılığı \\
overflow & taşma \\

%% --- P ---
page & sayfa \\
parity (RAID) & eşlik (RAID) \\
page fault & sayfa hatası \\
page table & sayfa tablosu \\
parallel computing & paralel hesaplama \\
parallel efficiency & paralel verimlilik \\
\texttt{PCI} Express & \texttt{PCI} Express \\
performance & başarım \\
pipeline & boru hattı \\
pipeline hazard & boru hattı tehlikesi \\
pipeline register & boru hattı kaydı \\
pipeline stage & boru hattı aşaması \\
port & bağlantı noktası \\
power consumption & güç tüketimi \\
power wall & güç duvarı \\
precise exception & kesin istisna \\
prefetching & önceden getirme \\
privileged mode & ayrıcalıklı kip \\
processor & işlemci \\
program counter & program sayacı \\
pseudo-instruction & sözde buyruk \\
pseudo-LRU & sözde-EUZK \\

%% --- Q ---
Quad-Level Cell (\texttt{QLC}) & \texttt{QLC} \\
quantum computing & kuantum hesaplama \\
qubit & kübit \\

%% --- R ---
R-type & R türü \\
race condition & yarış durumu \\
reconfigurable computing & yeniden yapılandırılabilir hesaplama \\
Redundant Array of Independent Disks (\texttt{RAID}) & \texttt{RAID} \\
Reduced Instruction Set Computer (\texttt{RISC}) & \texttt{RISC} \\
register & yazmaç \\
register file & yazmaç öbeği \\
register renaming & yazmaç yeniden adlandırma \\
Register Transfer Level (\texttt{RTL}) & yazmaç aktarım düzeyi (\texttt{RTL}) \\
relaxed consistency & gevşek tutarlılık \\
return address & dönüş adresi \\
ripple carry adder & dalga eldeli toplayıcı \\
\texttt{RISC-V} & \texttt{RISC-V} \\
roofline model & çatı çizgisi modeli \\
rotational latency & dönme gecikmesi \\
rounding & yuvarlama \\

%% --- S ---
S-type & S türü \\
\texttt{SATA} (Serial \texttt{ATA}) & \texttt{SATA} \\
secondary storage & ikincil bellek \\
seek time & arama zamanı \\
semaphore & semafor \\
sequential consistency & ardışık tutarlılık \\
Serial \texttt{ATA} (\texttt{SATA}) & \texttt{SATA} \\
set associative & küme ilişkili \\
shader & gölgelendirici \\
shared memory & paylaşımlı bellek \\
shift instructions & kaydırma buyrukları \\
side-channel attack & yan kanal saldırısı \\
sign bit & işaret biti \\
sign extension & işaretli genişletme \\
signed & işaretli \\
single-cycle processor & tek vuruşluk işlemci \\
Simultaneous Multithreading (\texttt{SMT}) & eşzamanlı çok kanallılık (\texttt{SMT}) \\
Single Instruction, Multiple Data (\texttt{SIMD}) & \texttt{SIMD} \\
Single Instruction, Multiple Threads (\texttt{SIMT}) & \texttt{SIMT} \\
Single-Level Cell (\texttt{SLC}) & \texttt{SLC} \\
single-cycle & tek çevrimli \\
snooping protocol & gözetleme protokolü \\
software & yazılım \\
scientific notation & bilimsel gösterim \\
significand / mantissa & anlamlı kısım \\
Solid State Drive (\texttt{SSD}) & katı hal sürücü (\texttt{SSD}) \\
spatial locality & uzamsal yerellik \\
\texttt{SPEC} & \texttt{SPEC} \\
Special Function Unit (\texttt{SFU}) & özel fonksiyon birimi (\texttt{SFU}) \\
Spectre & Spectre \\
speculative execution & spekülatif yürütme \\
spin-lock & döngüsel kilit \\
stack & yığıt \\
stack pointer & yığıt göstergesi \\
stale data & bayat veri \\
stall & durma \\
state register & durum yazmacı \\
static power consumption (leakage) & durağan güç tüketimi \\
static prediction & durağan öngörü \\
Static Random Access Memory (\texttt{SRAM}) & \texttt{SRAM} \\
store instruction & saklama buyruğu \\
Streaming Multiprocessor (\texttt{SM}) & Akış Çoklu İşlemci (\texttt{SM}) \\
Sv39 (39-bit virtual addressing) & Sv39 (RISC-V) \\
Sv48 (48-bit virtual addressing) & Sv48 (RISC-V) \\
Sv57 (57-bit virtual addressing) & Sv57 (RISC-V) \\
striping & şeritleme \\
structural hazard & yapısal tehlike \\
subroutine & altyordam \\
subtraction & çıkarma \\
superposition & süperpozisyon \\
superscalar & çoklu buyruk başlatımlı \\
synchronization & eşzamanlama \\
synthetic benchmarks & yapay sınama programları \\
System on Chip (\texttt{SoC}) & \texttt{SoC} \\
systolic array & sistolik dizi \\

%% --- T ---
tag & etiket \\
temporal locality & zamansal yerellik \\
three Cs model (3C) & üç C modeli (3C modeli) \\
Tensor Core & Tensör çekirdeği \\
Tensor Processing Unit (\texttt{TPU}) & Tensör İşlem Birimi (\texttt{TPU}) \\
texture memory & doku belleği \\
Thermal Design Power (\texttt{TDP}) & Soğutma Tasarım Gücü (\texttt{TDP}) \\
thread & iş parçacığı \\
thread pool & iş parçacığı havuzu \\
thread-level parallelism & iş parçacığı düzeyi paralelliği \\
throughput & verim / iş çıkarma oranı \\
Thunderbolt & Thunderbolt \\
training (AI) & eğitim (yapay zeka) \\
transistor & transistör \\
\texttt{TRIM} & \texttt{TRIM} \\
tournament predictor & turnuva öngörücüsü \\
Triple-Level Cell (\texttt{TLC}) & \texttt{TLC} \\
Translation Lookaside Buffer (\texttt{TLB}) & \texttt{TLB} \\
true data dependence & gerçek veri bağımlılığı \\
truth table & doğruluk tablosu \\
two's complement & ikiye tümleyen \\

%% --- U ---
U-type & U türü \\
unconditional branch & koşulsuz dallanma \\
Unicode & Unicode \\
Uniform Memory Access (\texttt{UMA}) & tekdüze bellek erişimi \\
unified memory architecture & birleşik bellek mimarisi \\
Universal Chiplet Interconnect Express (\texttt{UCIe}) & \texttt{UCIe} \\
Universal Serial Bus (\texttt{USB}) & \texttt{USB} \\
unsigned & işaretsiz \\
upper immediate instructions & üst anlık buyruklar \\

%% --- V ---
vacuum tube & vakum tüpü \\
Verilog & Verilog \\
vertical microprogramming & dikey mikroprogramlama \\
Very Large-Scale Integration (\texttt{VLSI}) & \texttt{VLSI} \\
\texttt{VHDL} & \texttt{VHDL} \\
virtual address & sanal adres \\
virtual memory & sanal bellek \\
Von Neumann architecture & Von Neumann mimarisi \\
Von Neumann bottleneck & Von Neumann darboğazı \\

%% --- W ---
warp & çözgü \\
weighted average & ağırlıklı ortalama \\
wear leveling & aşınma dengeleme \\
warp divergence & çözgü sapması \\
write-allocate & yazma tahsisli \\
write-back & geri yazma \\
write buffer & yazma ara belleği \\
write-invalidate & yazma geçersizleştirme \\
write propagation & yazma yayılımı \\
write serialization & yazma serileştirmesi \\
write-through & doğrudan yazma \\
write-update & yazma güncelleme \\

%% --- X ---
\texttt{XOR} gate & Özel VEYA kapısı \\

%% --- Z ---
zero extension & sıfırla genişletme \\

\end{longtable}
}
 ile dahil edilir.
% \chapter komutu ana-kitap.tex dosyasında tanımlanmalıdır.
% =============================================================

%% ============================================================
%%  BÖLÜM 1: Türkçe → İngilizce
%%  Türkçe alfabetik sıra: a,b,c,ç,d,e,f,g,ğ,h,ı,i,j,k,l,m,
%%                          n,o,ö,p,r,s,ş,t,u,ü,v,y,z
%% ============================================================
\section*{Türkçe -- İngilizce}
\addcontentsline{toc}{section}{Türkçe -- İngilizce}

{\small
\begin{longtable}{@{}p{7.6cm}p{7.6cm}@{}}
\toprule
\textbf{Türkçe Terim} & \textbf{İngilizce Karşılık} \\
\midrule
\endfirsthead

\toprule
\textbf{Türkçe Terim (devam)} & \textbf{İngilizce Karşılık} \\
\midrule
\endhead

\midrule
\multicolumn{2}{r}{\small\textit{(devamı sonraki sayfada)}} \\
\endfoot

\bottomrule
\endlastfoot

%% --- A ---
adresleme kipleri & addressing modes \\
ağırlıklı ortalama & weighted average \\
Akış Çoklu İşlemci (\texttt{SM}) & Streaming Multiprocessor (\texttt{SM}) \\
altyordam & subroutine \\
Amdahl Yasası & Amdahl's Law \\
\texttt{AMB} & \texttt{ALU} (Arithmetic Logic Unit) \\
anlamlı kısım & significand / mantissa \\
anlık değer & immediate value \\
anlık değer üreteci & immediate generator \\
ara bellek & buffer \\
ardışık tutarlılık & sequential consistency \\
arama zamanı & seek time \\
aritmetik buyruklar & arithmetic instructions \\
aşınma dengeleme & wear leveling \\
aritmetik mantık birimi (\texttt{AMB}) & Arithmetic Logic Unit (\texttt{ALU}) \\
aritmetik yoğunluk & arithmetic intensity \\
artık-N gösterimi & excess-N notation \\
\texttt{ASCII} & \texttt{ASCII} \\
\texttt{ASIC} & Application-Specific Integrated Circuit \\
çevirici dil (Assembly) & assembly language \\
atlama buyrukları & jump instructions \\
atomik işlem & atomic operation \\
aygıt & device \\
ayrıcalıklı kip & privileged mode \\

%% --- B ---
bilimsel gösterim & scientific notation \\
Booth algoritması & Booth's algorithm \\
B türü & B-type \\
bağlam değiştirme & context switching \\
bağlantı noktası & port \\
bağlayıcı & linker \\
bakış tablosu (\texttt{LUT}) & Look-Up Table (\texttt{LUT}) \\
bant genişliği & bandwidth \\
başarım & performance \\
bayat veri & stale data \\
bayt & byte \\
bayt sıralaması & byte order \\
büyük sayfa & huge page \\
\texttt{BBÇ} & \texttt{CPI} (cycles per instruction) \\
\texttt{BCD} & Binary-Coded Decimal \\
bellek & memory \\
bellek birleştirme & memory coalescing \\
bellek denetleyicisi & memory controller \\
bellek duvarı & memory wall \\
bellek koruması & memory protection \\
bellek hizalama & memory alignment \\
bellek hiyerarşisi & memory hierarchy \\
bellek tutarlığı & cache coherence \\
bellek tutarlık modeli & memory consistency model \\
bellekle eşlemeli G/Ç & memory-mapped I/O \\
benchmark & sınama programı (denektaşı) \\
big.LITTLE & big.LITTLE \\
bire tümleyen & one's complement \\
birleşik bellek mimarisi & unified memory architecture \\
birleşik devre & combinational circuit \\
bit & bit \\
blok & block \\
boru hattı & pipeline \\
boru hattı aşaması & pipeline stage \\
boru hattı kaydı & pipeline register \\
boru hattı tehlikesi & pipeline hazard \\
boşluk (boru hattı) & bubble / NOP \\
bölme & division \\
bulamama (önbellekte) & miss (cache miss) \\
bulamama cezası & miss penalty \\
bulamama oranı & miss rate \\
bulma (önbellekte) & hit (cache hit) \\
bulma oranı & hit rate \\
buyruk & instruction \\
buyruk başına çevrim (\texttt{BBÇ}) & cycles per instruction (\texttt{CPI}) \\
buyruk belleği & instruction memory \\
buyruk çözme & instruction decode \\
buyruk düzeyi paralelliği & instruction-level parallelism (\texttt{ILP}) \\
buyruk biçimi & instruction format \\
buyruk getirme birimi & instruction fetch unit \\
buyruk kümesi mimarisi & Instruction Set Architecture (\texttt{ISA}) \\
buyruk önbelleği & instruction cache \\
buyruk sayacı & instruction counter \\
büyüğü başta & big-endian \\

%% --- C ---
chiplet & chiplet \\
\texttt{CISC} & Complex Instruction Set Computer \\
\texttt{CMOS} & Complementary \texttt{MOS} \\
\texttt{CUDA} & Compute Unified Device Architecture \\

%% --- Ç ---
\texttt{ÇBB} & \texttt{IPC} (instructions per cycle) \\
çağrı kuralları & calling conventions \\
çarpma & multiplication \\
çarpma-toplama (\texttt{MAC}) & multiply-accumulate (\texttt{MAC}) \\
çatı çizgisi modeli & roofline model \\
çekirdek & core \\
çekirdek fonksiyonu (GPU) & kernel function \\
çerçeve göstergesi & frame pointer \\
çevirici & assembler \\
çevrim & cycle \\
çevrim başına buyruk (\texttt{ÇBB}) & instructions per cycle (\texttt{IPC}) \\
çevrim zamanı & cycle time \\
çıkarım (yapay zeka) & inference \\
çıkarma & subtraction \\
yonga içi bağlantı ağı (\texttt{NoC}) & Network on Chip (\texttt{NoC}) \\
çok çekirdekli işlemci & multi-core processor \\
çok seviyeli sayfa tablosu & multi-level page table \\
çok çevrimli & multi-cycle \\
çok vuruşluk işlemci & multi-cycle processor \\
çoklayıcı & multiplexer \\
çözgü & warp \\
çözgü sapması & warp divergence \\

%% --- D ---
dağıtık bellek & distributed memory \\
durum yazmacı & state register \\
dalga eldeli toplayıcı & ripple carry adder \\
dallanma buyrukları & branch instructions \\
dallanma hedef ara belleği & branch target buffer (BTB) \\
dallanma öngörüsü & branch prediction \\
darboğaz & bottleneck \\
\texttt{DDR} & Double Data Rate \\
DEĞİL kapısı & \texttt{NOT} gate \\
dekoder & decoder \\
denetim birimi & control unit \\
denetim sinyali & control signal \\
denetim tablosu & control table \\
denetim tehlikesi & control hazard \\
devingen güç tüketimi & dynamic power consumption \\
devingen öngörü & dynamic prediction \\
Dennard ölçeklemesi & Dennard scaling \\
derin öğrenme & deep learning \\
derleyici & compiler \\
dikey mikroprogramlama & vertical microprogramming \\
dizin tabanlı protokol & directory-based protocol \\
\texttt{DMA} & Direct Memory Access \\
\texttt{DRAM} & Dynamic Random Access Memory \\
doğrudan eşlemeli önbellek & direct-mapped cache \\
doğruluk tablosu & truth table \\
dönme gecikmesi & rotational latency \\
doğrudan yazma & write-through \\
doku belleği & texture memory \\
dolanıklık & entanglement \\
doluluk & occupancy \\
donanım & hardware \\
donanımsal çok kanallılık & hardware multithreading \\
dönüş adresi & return address \\
döngüsel kilit & spin-lock \\
durağan öngörü & static prediction \\
durağan güç tüketimi & static power consumption (leakage) \\
durma & stall \\
düzenleme birimi & hazard detection unit \\

%% --- E ---
\texttt{EDVAC} & \texttt{EDVAC} \\
eğitim (yapay zeka) & training \\
elde & carry \\
elden önceden üretme & carry lookahead \\
\texttt{ENIAC} & \texttt{ENIAC} \\
engel & barrier \\
eşlik (RAID) & parity (RAID) \\
erişim süresi & access time \\
eşzamanlama & synchronization \\
eşzamanlı çok kanallılık (\texttt{SMT}) & Simultaneous Multithreading (\texttt{SMT}) \\
eşzamanlılık & concurrency \\
etiket & tag \\
etkin adres & effective address \\

%% --- F ---
Flaş çeviri katmanı (\texttt{FTL}) & Flash Translation Layer (\texttt{FTL}) \\
flaş bellek & flash memory \\
flip-flop & flip-flop \\
\texttt{FPGA} & Field-Programmable Gate Array \\

%% --- G ---
gecikme & latency \\
geometrik ortalama & geometric mean \\
gecikme odaklı tasarım & latency-oriented design \\
geciktirilmiş dallanma & delayed branching \\
gerçek veri bağımlılığı & true data dependence \\
genel amaçlı yazmaç & general-purpose register \\
genel bellek & global memory \\
geri yazma & write-back \\
getir-çöz-yürüt döngüsü & fetch-decode-execute cycle \\
gevşek tutarlılık & relaxed consistency \\
giriş/çıkış (\texttt{G/Ç}) & input/output (\texttt{I/O}) \\
gölgelendirici & shader \\
gözetleme protokolü & snooping protocol \\
\texttt{GPGPU} & General-Purpose computing on GPU \\
\texttt{GPU} & Graphics Processing Unit \\
gshare öngörücü & gshare predictor \\
Gustafson Yasası & Gustafson's Law \\
güç duvarı & power wall \\
güç tüketimi & power consumption \\

%% --- Ğ ---

%% --- H ---
Harvard mimarisi & Harvard architecture \\
\texttt{HBM} & High Bandwidth Memory \\
heterojen hesaplama & heterogeneous computing \\

%% --- Önbellekte bulamama terimleri (B harfi altında da var) ---
%% --- I ---
ızgara & grid \\

%% --- İ ---
I türü & I-type \\
\texttt{IEEE 754} & \texttt{IEEE 754} \\
yol & bus \\
ikincil bellek & secondary storage \\
\texttt{IOPS} & Input/Output Operations Per Second \\
ikiye tümleyen & two's complement \\
ikilik taban & binary \\
%% isabet → bulma olarak B harfi altında (bkz. "bulma")
iş çıkarma oranı & throughput \\
iş parçacığı & thread \\
iş parçacığı düzeyi paralelliği & thread-level parallelism \\
iş parçacığı havuzu & thread pool \\
işaret biti & sign bit \\
işaretli genişletme & sign extension \\
işaretli & signed \\
işaretsiz & unsigned \\
işlem kodu & opcode \\
işlemci & processor \\
işletim sistemi & operating system \\
istisna (boru hattı) & exception (pipeline) \\

%% --- J ---
J türü & J-type \\

%% --- K ---
kritik yol & critical path \\
karşı bağımlılık & anti-dependence \\
karşılaştırma buyrukları & comparison instructions \\
kaydırma buyrukları & shift instructions \\
kayan nokta & floating point \\
kayan noktalı sayı birimi & floating point unit (\texttt{FPU}) \\
katı hal sürücü (\texttt{SSD}) & Solid State Drive (\texttt{SSD}) \\
kesin istisna & precise exception \\
kesme & interrupt \\
kesme denetleyicisi & interrupt controller \\
kesme hizmet yordamı & interrupt service routine \\
kilit & lock/mutex \\
kilitlenme & deadlock \\
kod çözücü & decoder \\
koşullu dallanma & conditional branch \\
koşulsuz dallanma & unconditional branch \\
kuantum hesaplama & quantum computing \\
kübit & qubit \\
küçüğü başta & little-endian \\
kırmık & die \\
küme ilişkili & set associative \\

%% --- L ---
L1 önbellek & L1 cache \\
L2 önbellek & L2 cache \\
L3 önbellek & L3 cache \\

%% --- M ---
\texttt{M.2} & \texttt{M.2} (form factor) \\
makine dili & machine language \\
makro buyruk & macro instruction \\
\texttt{MLC} & Multi-Level Cell \\
mantık buyrukları & logic instructions \\
mantık kapısı & logic gate \\
matris çarpımı & matrix multiplication \\
Meltdown & Meltdown \\
\texttt{MESI} protokolü & \texttt{MESI} protocol \\
merkezi işlem birimi (\texttt{MİB}) & Central Processing Unit (\texttt{CPU}) \\
mesaj geçirme & message passing \\
\texttt{MFLOPS} & Millions of Floating Point Operations Per Second \\
mikro buyruk & micro-operation \\
mikro program belleği & control store \\
mikromimari & microarchitecture \\
mikroprogramlama & microprogramming \\
\texttt{MIPS} & Millions of Instructions Per Second \\
\texttt{MİB} & \texttt{CPU} (Central Processing Unit) \\
Moore Yasası & Moore's Law \\

%% --- N ---
\texttt{NAND} Flaş bellek & \texttt{NAND Flash} memory \\
\texttt{NUMA} & Non-Uniform Memory Access \\
\texttt{NVMe} & Non-Volatile Memory Express \\

%% --- O ---
olağanlaştırma & normalization \\
on altılık taban & hexadecimal \\
optimal çıkarma algoritması & optimal (Bélády's) replacement algorithm \\
ortalama bellek erişim süresi (\texttt{AMAT}) & Average Memory Access Time (\texttt{AMAT}) \\

%% --- Ö ---
önbellek & cache \\
önbellek satırı & cache line \\
önceden getirme & prefetching \\
Özel VEYA kapısı & \texttt{XOR} gate \\
özel fonksiyon birimi (\texttt{SFU}) & Special Function Unit (\texttt{SFU}) \\

%% --- P ---
paralel hesaplama & parallel computing \\
paralel verimlilik & parallel efficiency \\
paylaşımlı bellek & shared memory \\
\texttt{PCI} Express & \texttt{PCI} Express \\
program sayacı & program counter \\

%% --- Q ---
\texttt{QLC} & Quad-Level Cell \\

%% --- R ---
R türü & R-type \\
\texttt{RAID} & Redundant Array of Independent Disks \\
\texttt{RISC} & Reduced Instruction Set Computer \\
\texttt{RISC-V} & \texttt{RISC-V} \\

%% --- S ---
S türü & S-type \\
sıfırla genişletme & zero extension \\
saat çevrimi & clock cycle \\
saat sıklığı & clock frequency \\
sabit bellek & constant memory \\
sabit disk sürücüsü (\texttt{HDD}) & hard disk drive (\texttt{HDD}) \\
\texttt{SATA} & Serial \texttt{ATA} \\
saklama buyruğu & store instruction \\
sanal adres & virtual address \\
sanal bellek & virtual memory \\
sayfa & page \\
sayfa hatası & page fault \\
sayfa tablosu & page table \\
sekizlik taban & octal \\
semafor & semaphore \\
sıra dışı yürütme & out-of-order execution \\
şeritleme & striping \\
\texttt{SIMD} & Single Instruction, Multiple Data \\
\texttt{SLC} & Single-Level Cell \\
\texttt{SIMT} & Single Instruction, Multiple Threads \\
Sinir Ağı İşlem Birimi (\texttt{NPU}) & Neural Processing Unit (\texttt{NPU}) \\
sınama programı (denektaşı) & benchmark \\
sistolik dizi & systolic array \\
Soğutma Tasarım Gücü (\texttt{TDP}) & Thermal Design Power (\texttt{TDP}) \\
\texttt{SoC} & System on Chip \\
sonlu durum makinesi & finite state machine \\
sonuç bağımlılığı & output dependence \\
sonradan yazma & write-back \\
soyutlama & abstraction \\
sözde buyruk & pseudo-instruction \\
sözde-EUZK & pseudo-LRU \\
Spectre & Spectre \\
\texttt{SPEC} & \texttt{SPEC} \\
spekülatif yürütme & speculative execution \\
\texttt{SRAM} & Static Random Access Memory \\
Sv39 (RISC-V) & Sv39 (39-bit virtual addressing) \\
Sv48 (RISC-V) & Sv48 (48-bit virtual addressing) \\
Sv57 (RISC-V) & Sv57 (57-bit virtual addressing) \\
çoklu buyruk başlatımlı & superscalar \\
süperpozisyon & superposition \\

%% --- Ş ---

%% --- T ---
tam ilişkili & fully associative \\
Thunderbolt & Thunderbolt \\
\texttt{TLC} & Triple-Level Cell \\
\texttt{TRIM} & \texttt{TRIM} \\
tam sayı & integer \\
tam toplayıcı & full adder \\
taşma & overflow \\
tek çevrimli & single-cycle \\
tek vuruşluk işlemci & single-cycle processor \\
tekdüze bellek erişimi & Uniform Memory Access (\texttt{UMA}) \\
tekdüze olmayan bellek erişimi & Non-Uniform Memory Access (\texttt{NUMA}) \\
Tensör çekirdeği & Tensor Core \\
Tensör İşlem Birimi (\texttt{TPU}) & Tensor Processing Unit (\texttt{TPU}) \\
\texttt{TLB} & Translation Lookaside Buffer \\
toplama & addition \\
transistör & transistor \\
turnuva öngörücüsü & tournament predictor \\
tümleşik devre & integrated circuit \\

%% --- U ---
uç bilgi işlem & edge computing \\
U türü & U-type \\
\texttt{USB} & Universal Serial Bus \\
\texttt{UCIe} & Universal Chiplet Interconnect Express \\
\texttt{UMA} & Uniform Memory Access \\
Unicode & Unicode \\
uzamsal yerellik & spatial locality \\

%% --- Ü ---
üç C modeli (3C modeli) & three Cs model \\
üs & exponent \\
üst anlık buyruklar & upper immediate instructions \\

%% --- V ---
vakum tüpü & vacuum tube \\
VE kapısı & \texttt{AND} gate \\
veri akış mimarisi & dataflow architecture \\
veri aktarma buyrukları & data transfer instructions \\
veri bağımlılığı & data dependency \\
veri belleği & data memory \\
veri tehlikesi & data hazard \\
veri yolu & datapath \\
veri yönlendirmesi & forwarding \\
verim & throughput \\
verim odaklı tasarım & throughput-oriented design \\
Verilog & Verilog \\
VEYA kapısı & \texttt{OR} gate \\
\texttt{VHDL} & \texttt{VHDL} \\
\texttt{VLSI} & Very Large-Scale Integration \\
Von Neumann darboğazı & Von Neumann bottleneck \\
Von Neumann mimarisi & Von Neumann architecture \\

%% --- Y ---
yan kanal saldırısı & side-channel attack \\
yanlış paylaşım & false sharing \\
yapay sınama programları & synthetic benchmarks \\
yapay zeka & artificial intelligence \\
yapısal tehlike & structural hazard \\
yarım toplayıcı & half adder \\
yarış durumu & race condition \\
yatay mikroprogramlama & horizontal microprogramming \\
yazılım & software \\
yazma geçersizleştirme & write-invalidate \\
yazma güncelleme & write-update \\
yazma serileştirmesi & write serialization \\
yazma tahsisli & write-allocate \\
yazma tahsissiz & no-write-allocate \\
yazma ara belleği & write buffer \\
yazma yayılımı & write propagation \\
yazmaç & register \\
yazmaç aktarım düzeyi (\texttt{RTL}) & Register Transfer Level (\texttt{RTL}) \\
yazmaç öbeği & register file \\
yazmaç yeniden adlandırma & register renaming \\
yekpare kırmık & monolithic die \\
yeniden yapılandırılabilir hesaplama & reconfigurable computing \\
yerel bellek & local memory \\
yerellik ilkesi & locality principle \\
yığıt & stack \\
yığıt göstergesi & stack pointer \\
yonga & chip \\
yuvarlama & rounding \\
yükleyici & loader \\
yükle-kullan sorunu & load-use hazard \\
yükleme buyruğu & load instruction \\
yürütme birimi & execution unit \\
yürütme zamanı & execution time \\

%% --- Z ---
zamansal yerellik & temporal locality \\

\end{longtable}
}


\bigskip

%% ============================================================
%%  BÖLÜM 2: İngilizce → Türkçe
%%  İngilizce alfabetik sıra (A-Z)
%% ============================================================
\section*{İngilizce -- Türkçe}
\addcontentsline{toc}{section}{İngilizce -- Türkçe}

{\small
\begin{longtable}{@{}p{7.6cm}p{7.6cm}@{}}
\toprule
\textbf{İngilizce Terim} & \textbf{Türkçe Karşılık} \\
\midrule
\endfirsthead

\toprule
\textbf{İngilizce Terim (cont.)} & \textbf{Türkçe Karşılık} \\
\midrule
\endhead

\midrule
\multicolumn{2}{r}{\small\textit{(continued on next page)}} \\
\endfoot

\bottomrule
\endlastfoot

%% --- A ---
abstraction & soyutlama \\
access time & erişim süresi \\
addition & toplama \\
Average Memory Access Time (\texttt{AMAT}) & ortalama bellek erişim süresi (\texttt{AMAT}) \\
addressing modes & adresleme kipleri \\
Amdahl's Law & Amdahl Yasası \\
\texttt{AND} gate & VE kapısı \\
anti-dependence & karşı bağımlılık \\
Application-Specific Integrated Circuit (\texttt{ASIC}) & \texttt{ASIC} \\
Arithmetic Logic Unit (\texttt{ALU}) & aritmetik mantık birimi (\texttt{AMB}) \\
arithmetic instructions & aritmetik buyruklar \\
arithmetic intensity & aritmetik yoğunluk \\
artificial intelligence & yapay zeka \\
\texttt{ASCII} & \texttt{ASCII} \\
assembler & çevirici \\
assembly language & çevirici dil \\
atomic operation & atomik işlem \\

%% --- B ---
B-type & B türü \\
Booth's algorithm & Booth algoritması \\
bandwidth & bant genişliği \\
barrier & engel \\
benchmark & sınama programı (denektaşı) \\
big-endian & büyüğü başta \\
big.LITTLE & big.LITTLE \\
binary & ikilik taban \\
Binary-Coded Decimal (\texttt{BCD}) & \texttt{BCD} \\
bit & bit \\
block & blok \\
bottleneck & darboğaz \\
branch instructions & dallanma buyrukları \\
branch prediction & dallanma öngörüsü \\
branch target buffer (BTB) & dallanma hedef ara belleği \\
buffer & ara bellek \\
bubble (pipeline) & boşluk (boru hattı) \\
bus & yol \\
byte & bayt \\
byte order & bayt sıralaması \\

%% --- C ---
cache & önbellek \\
cache coherence & bellek tutarlığı \\
cache line & önbellek satırı \\
calling conventions & çağrı kuralları \\
carry & elde \\
carry lookahead & elden önceden üretme \\
Central Processing Unit (\texttt{CPU}) & merkezi işlem birimi (\texttt{MİB}) \\
chip & yonga \\
chiplet & chiplet \\
clock cycle & saat çevrimi \\
clock frequency & saat sıklığı \\
combinational circuit & birleşik devre \\
Complementary \texttt{MOS} (\texttt{CMOS}) & \texttt{CMOS} \\
compiler & derleyici \\
comparison instructions & karşılaştırma buyrukları \\
Complex Instruction Set Computer (\texttt{CISC}) & \texttt{CISC} \\
Compute Unified Device Architecture (\texttt{CUDA}) & \texttt{CUDA} \\
concurrency & eşzamanlılık \\
conditional branch & koşullu dallanma \\
constant memory & sabit bellek \\
context switching & bağlam değiştirme \\
control hazard & denetim tehlikesi \\
control signal & denetim sinyali \\
control store & mikro program belleği \\
control table & denetim tablosu \\
control unit & denetim birimi \\
core & çekirdek \\
critical path & kritik yol \\
cycle & çevrim \\
cycle time & çevrim zamanı \\
cycles per instruction (\texttt{CPI}) & buyruk başına çevrim (\texttt{BBÇ}) \\

%% --- D ---
die & kırmık \\
data dependency & veri bağımlılığı \\
data hazard & veri tehlikesi \\
data memory & veri belleği \\
data transfer instructions & veri aktarma buyrukları \\
dataflow architecture & veri akış mimarisi \\
datapath & veri yolu \\
deadlock & kilitlenme \\
decoder & dekoder / kod çözücü \\
deep learning & derin öğrenme \\
delayed branching & geciktirilmiş dallanma \\
Dennard scaling & Dennard ölçeklemesi \\
device & aygıt \\
Direct Memory Access (\texttt{DMA}) & \texttt{DMA} \\
direct-mapped cache & doğrudan eşlemeli önbellek \\
directory-based protocol & dizin tabanlı protokol \\
distributed memory & dağıtık bellek \\
division & bölme \\
Double Data Rate (\texttt{DDR}) & \texttt{DDR} \\
dynamic prediction & devingen öngörü \\
dynamic power consumption & devingen güç tüketimi \\
Dynamic Random Access Memory (\texttt{DRAM}) & \texttt{DRAM} \\

%% --- E ---
edge computing & uç bilgi işlem \\
\texttt{EDVAC} & \texttt{EDVAC} \\
effective address & etkin adres \\
\texttt{ENIAC} & \texttt{ENIAC} \\
entanglement & dolanıklık \\
exception (pipeline) & istisna (boru hattı) \\
excess-N notation & artık-N gösterimi \\
execution time & yürütme zamanı \\
execution unit & yürütme birimi \\
exponent & üs \\

%% --- F ---
false sharing & yanlış paylaşım \\
fetch-decode-execute cycle & getir-çöz-yürüt döngüsü \\
Field-Programmable Gate Array (\texttt{FPGA}) & \texttt{FPGA} \\
finite state machine & sonlu durum makinesi \\
flash memory & flaş bellek \\
Flash Translation Layer (\texttt{FTL}) & Flaş çeviri katmanı (\texttt{FTL}) \\
flip-flop & flip-flop \\
floating point & kayan nokta \\
floating point unit (\texttt{FPU}) & kayan noktalı sayı birimi \\
forwarding & veri yönlendirmesi \\
frame pointer & çerçeve göstergesi \\
full adder & tam toplayıcı \\
fully associative & tam ilişkili \\

%% --- G ---
General-Purpose computing on GPU (\texttt{GPGPU}) & \texttt{GPGPU} \\
general-purpose register & genel amaçlı yazmaç \\
geometric mean & geometrik ortalama \\
global memory & genel bellek \\
Graphics Processing Unit (\texttt{GPU}) & \texttt{GPU} \\
grid & ızgara \\
gshare predictor & gshare öngörücü \\
Gustafson's Law & Gustafson Yasası \\

%% --- H ---
half adder & yarım toplayıcı \\
hard disk drive (\texttt{HDD}) & sabit disk sürücüsü \\
hardware & donanım \\
hardware multithreading & donanımsal çok kanallılık \\
Harvard architecture & Harvard mimarisi \\
hazard detection unit & düzenleme birimi \\
heterogeneous computing & heterojen hesaplama \\
hexadecimal & on altılık taban \\
High Bandwidth Memory (\texttt{HBM}) & \texttt{HBM} \\
hit (cache hit) & bulma (önbellekte) \\
hit rate & bulma oranı \\
horizontal microprogramming & yatay mikroprogramlama \\
huge page & büyük sayfa \\

%% --- I ---
I-type & I türü \\
\texttt{IEEE 754} & \texttt{IEEE 754} \\
immediate generator & anlık değer üreteci \\
immediate value & anlık değer \\
inference (AI) & çıkarım (yapay zeka) \\
Input/Output Operations Per Second (\texttt{IOPS}) & \texttt{IOPS} \\
input/output (\texttt{I/O}) & giriş/çıkış (\texttt{G/Ç}) \\
instruction & buyruk \\
instruction cache & buyruk önbelleği \\
instruction decode & buyruk çözme \\
instruction counter & buyruk sayacı \\
instruction fetch unit & buyruk getirme birimi \\
instruction format & buyruk biçimi \\
instruction memory & buyruk belleği \\
Instruction Set Architecture (\texttt{ISA}) & buyruk kümesi mimarisi \\
instruction-level parallelism (\texttt{ILP}) & buyruk düzeyi paralelliği \\
instructions per cycle (\texttt{IPC}) & çevrim başına buyruk (\texttt{ÇBB}) \\
integer & tam sayı \\
integrated circuit & tümleşik devre \\
interrupt & kesme \\
interrupt controller & kesme denetleyicisi \\
interrupt service routine & kesme hizmet yordamı \\

%% --- J ---
J-type & J türü \\
jump instructions & atlama buyrukları \\

%% --- K ---
kernel function (GPU) & çekirdek fonksiyonu (GPU) \\

%% --- L ---
L1 cache & L1 önbellek \\
L2 cache & L2 önbellek \\
L3 cache & L3 önbellek \\
latency & gecikme \\
latency-oriented design & gecikme odaklı tasarım \\
linker & bağlayıcı \\
little-endian & küçüğü başta \\
load instruction & yükleme buyruğu \\
load-use hazard & yükle-kullan sorunu \\
loader & yükleyici \\
local memory & yerel bellek \\
locality principle & yerellik ilkesi \\
Look-Up Table (\texttt{LUT}) & bakış tablosu (\texttt{LUT}) \\
lock/mutex & kilit \\
logic gate & mantık kapısı \\
logic instructions & mantık buyrukları \\

%% --- M ---
\texttt{M.2} & \texttt{M.2} (form faktörü) \\
machine language & makine dili \\
macro instruction & makro buyruk \\
matrix multiplication & matris çarpımı \\
Meltdown & Meltdown \\
memory & bellek \\
memory alignment & bellek hizalama \\
memory coalescing & bellek birleştirme \\
memory consistency model & bellek tutarlık modeli \\
memory controller & bellek denetleyicisi \\
memory hierarchy & bellek hiyerarşisi \\
memory protection & bellek koruması \\
memory wall & bellek duvarı \\
memory-mapped I/O & bellekle eşlemeli G/Ç \\
\texttt{MESI} protocol & \texttt{MESI} protokolü \\
message passing & mesaj geçirme \\
microarchitecture & mikromimari \\
micro-operation & mikro buyruk \\
microprogramming & mikroprogramlama \\
Millions of Floating Point Operations Per Second (\texttt{MFLOPS}) & \texttt{MFLOPS} \\
Millions of Instructions Per Second (\texttt{MIPS}) & \texttt{MIPS} \\
miss (cache miss) & bulamama (önbellekte) \\
miss penalty & bulamama cezası \\
miss rate & bulamama oranı \\
monolithic die & yekpare kırmık \\
Moore's Law & Moore Yasası \\
multi-core processor & çok çekirdekli işlemci \\
multiply-accumulate (\texttt{MAC}) & çarpma-toplama (\texttt{MAC}) \\
multi-cycle & çok çevrimli \\
multi-cycle processor & çok vuruşluk işlemci \\
Multi-Level Cell (\texttt{MLC}) & \texttt{MLC} \\
multi-level page table & çok seviyeli sayfa tablosu \\
multiplication & çarpma \\
multiplexer & çoklayıcı \\

%% --- N ---
\texttt{NAND Flash} memory & \texttt{NAND} Flaş bellek \\
Network on Chip (\texttt{NoC}) & yonga içi bağlantı ağı (\texttt{NoC}) \\
Neural Processing Unit (\texttt{NPU}) & Sinir Ağı İşlem Birimi (\texttt{NPU}) \\
Non-Uniform Memory Access (\texttt{NUMA}) & tekdüze olmayan bellek erişimi \\
Non-Volatile Memory Express (\texttt{NVMe}) & \texttt{NVMe} \\
no-write-allocate & yazma tahsissiz \\
normalization & olağanlaştırma \\
\texttt{NOT} gate & DEĞİL kapısı \\

%% --- O ---
occupancy & doluluk \\
octal & sekizlik taban \\
one's complement & bire tümleyen \\
opcode & işlem kodu \\
operating system & işletim sistemi \\
optimal (Bélády's) replacement algorithm & optimal çıkarma algoritması \\
\texttt{OR} gate & VEYA kapısı \\
out-of-order execution & sıra dışı yürütme \\
output dependence & sonuç bağımlılığı \\
overflow & taşma \\

%% --- P ---
page & sayfa \\
parity (RAID) & eşlik (RAID) \\
page fault & sayfa hatası \\
page table & sayfa tablosu \\
parallel computing & paralel hesaplama \\
parallel efficiency & paralel verimlilik \\
\texttt{PCI} Express & \texttt{PCI} Express \\
performance & başarım \\
pipeline & boru hattı \\
pipeline hazard & boru hattı tehlikesi \\
pipeline register & boru hattı kaydı \\
pipeline stage & boru hattı aşaması \\
port & bağlantı noktası \\
power consumption & güç tüketimi \\
power wall & güç duvarı \\
precise exception & kesin istisna \\
prefetching & önceden getirme \\
privileged mode & ayrıcalıklı kip \\
processor & işlemci \\
program counter & program sayacı \\
pseudo-instruction & sözde buyruk \\
pseudo-LRU & sözde-EUZK \\

%% --- Q ---
Quad-Level Cell (\texttt{QLC}) & \texttt{QLC} \\
quantum computing & kuantum hesaplama \\
qubit & kübit \\

%% --- R ---
R-type & R türü \\
race condition & yarış durumu \\
reconfigurable computing & yeniden yapılandırılabilir hesaplama \\
Redundant Array of Independent Disks (\texttt{RAID}) & \texttt{RAID} \\
Reduced Instruction Set Computer (\texttt{RISC}) & \texttt{RISC} \\
register & yazmaç \\
register file & yazmaç öbeği \\
register renaming & yazmaç yeniden adlandırma \\
Register Transfer Level (\texttt{RTL}) & yazmaç aktarım düzeyi (\texttt{RTL}) \\
relaxed consistency & gevşek tutarlılık \\
return address & dönüş adresi \\
ripple carry adder & dalga eldeli toplayıcı \\
\texttt{RISC-V} & \texttt{RISC-V} \\
roofline model & çatı çizgisi modeli \\
rotational latency & dönme gecikmesi \\
rounding & yuvarlama \\

%% --- S ---
S-type & S türü \\
\texttt{SATA} (Serial \texttt{ATA}) & \texttt{SATA} \\
secondary storage & ikincil bellek \\
seek time & arama zamanı \\
semaphore & semafor \\
sequential consistency & ardışık tutarlılık \\
Serial \texttt{ATA} (\texttt{SATA}) & \texttt{SATA} \\
set associative & küme ilişkili \\
shader & gölgelendirici \\
shared memory & paylaşımlı bellek \\
shift instructions & kaydırma buyrukları \\
side-channel attack & yan kanal saldırısı \\
sign bit & işaret biti \\
sign extension & işaretli genişletme \\
signed & işaretli \\
single-cycle processor & tek vuruşluk işlemci \\
Simultaneous Multithreading (\texttt{SMT}) & eşzamanlı çok kanallılık (\texttt{SMT}) \\
Single Instruction, Multiple Data (\texttt{SIMD}) & \texttt{SIMD} \\
Single Instruction, Multiple Threads (\texttt{SIMT}) & \texttt{SIMT} \\
Single-Level Cell (\texttt{SLC}) & \texttt{SLC} \\
single-cycle & tek çevrimli \\
snooping protocol & gözetleme protokolü \\
software & yazılım \\
scientific notation & bilimsel gösterim \\
significand / mantissa & anlamlı kısım \\
Solid State Drive (\texttt{SSD}) & katı hal sürücü (\texttt{SSD}) \\
spatial locality & uzamsal yerellik \\
\texttt{SPEC} & \texttt{SPEC} \\
Special Function Unit (\texttt{SFU}) & özel fonksiyon birimi (\texttt{SFU}) \\
Spectre & Spectre \\
speculative execution & spekülatif yürütme \\
spin-lock & döngüsel kilit \\
stack & yığıt \\
stack pointer & yığıt göstergesi \\
stale data & bayat veri \\
stall & durma \\
state register & durum yazmacı \\
static power consumption (leakage) & durağan güç tüketimi \\
static prediction & durağan öngörü \\
Static Random Access Memory (\texttt{SRAM}) & \texttt{SRAM} \\
store instruction & saklama buyruğu \\
Streaming Multiprocessor (\texttt{SM}) & Akış Çoklu İşlemci (\texttt{SM}) \\
Sv39 (39-bit virtual addressing) & Sv39 (RISC-V) \\
Sv48 (48-bit virtual addressing) & Sv48 (RISC-V) \\
Sv57 (57-bit virtual addressing) & Sv57 (RISC-V) \\
striping & şeritleme \\
structural hazard & yapısal tehlike \\
subroutine & altyordam \\
subtraction & çıkarma \\
superposition & süperpozisyon \\
superscalar & çoklu buyruk başlatımlı \\
synchronization & eşzamanlama \\
synthetic benchmarks & yapay sınama programları \\
System on Chip (\texttt{SoC}) & \texttt{SoC} \\
systolic array & sistolik dizi \\

%% --- T ---
tag & etiket \\
temporal locality & zamansal yerellik \\
three Cs model (3C) & üç C modeli (3C modeli) \\
Tensor Core & Tensör çekirdeği \\
Tensor Processing Unit (\texttt{TPU}) & Tensör İşlem Birimi (\texttt{TPU}) \\
texture memory & doku belleği \\
Thermal Design Power (\texttt{TDP}) & Soğutma Tasarım Gücü (\texttt{TDP}) \\
thread & iş parçacığı \\
thread pool & iş parçacığı havuzu \\
thread-level parallelism & iş parçacığı düzeyi paralelliği \\
throughput & verim / iş çıkarma oranı \\
Thunderbolt & Thunderbolt \\
training (AI) & eğitim (yapay zeka) \\
transistor & transistör \\
\texttt{TRIM} & \texttt{TRIM} \\
tournament predictor & turnuva öngörücüsü \\
Triple-Level Cell (\texttt{TLC}) & \texttt{TLC} \\
Translation Lookaside Buffer (\texttt{TLB}) & \texttt{TLB} \\
true data dependence & gerçek veri bağımlılığı \\
truth table & doğruluk tablosu \\
two's complement & ikiye tümleyen \\

%% --- U ---
U-type & U türü \\
unconditional branch & koşulsuz dallanma \\
Unicode & Unicode \\
Uniform Memory Access (\texttt{UMA}) & tekdüze bellek erişimi \\
unified memory architecture & birleşik bellek mimarisi \\
Universal Chiplet Interconnect Express (\texttt{UCIe}) & \texttt{UCIe} \\
Universal Serial Bus (\texttt{USB}) & \texttt{USB} \\
unsigned & işaretsiz \\
upper immediate instructions & üst anlık buyruklar \\

%% --- V ---
vacuum tube & vakum tüpü \\
Verilog & Verilog \\
vertical microprogramming & dikey mikroprogramlama \\
Very Large-Scale Integration (\texttt{VLSI}) & \texttt{VLSI} \\
\texttt{VHDL} & \texttt{VHDL} \\
virtual address & sanal adres \\
virtual memory & sanal bellek \\
Von Neumann architecture & Von Neumann mimarisi \\
Von Neumann bottleneck & Von Neumann darboğazı \\

%% --- W ---
warp & çözgü \\
weighted average & ağırlıklı ortalama \\
wear leveling & aşınma dengeleme \\
warp divergence & çözgü sapması \\
write-allocate & yazma tahsisli \\
write-back & geri yazma \\
write buffer & yazma ara belleği \\
write-invalidate & yazma geçersizleştirme \\
write propagation & yazma yayılımı \\
write serialization & yazma serileştirmesi \\
write-through & doğrudan yazma \\
write-update & yazma güncelleme \\

%% --- X ---
\texttt{XOR} gate & Özel VEYA kapısı \\

%% --- Z ---
zero extension & sıfırla genişletme \\

\end{longtable}
}
